\documentclass[11pt]{amsart}

% Encoding and fonts
\usepackage[T1]{fontenc}
\usepackage[utf8]{inputenc}
\usepackage{lmodern}

% Math packages
\usepackage{amsmath}
\usepackage{amssymb}
\usepackage{amsfonts}
\usepackage{amsthm}
\usepackage{mathtools}

% Layout and utilities
\usepackage{geometry}
\usepackage{enumitem}
\usepackage[hidelinks]{hyperref}

\geometry{margin=1in}

\numberwithin{equation}{section}
\setcounter{tocdepth}{3}
\setcounter{secnumdepth}{3}

% ---------------------------------------------------------
% Theorem-like environments
% ---------------------------------------------------------

\theoremstyle{plain}
\newtheorem{theorem}{Theorem}[section]
\newtheorem{proposition}[theorem]{Proposition}
\newtheorem{lemma}[theorem]{Lemma}
\newtheorem{corollary}[theorem]{Corollary}
\newtheorem{conjecture}[theorem]{Conjecture}

\theoremstyle{definition}
\newtheorem{definition}[theorem]{Definition}
\newtheorem{heuristic}[theorem]{Heuristic}
\newtheorem{question}[theorem]{Question}

\theoremstyle{remark}
\newtheorem{remark}[theorem]{Remark}
\newtheorem{warning}[theorem]{Warning}

\title{Exponential Sums over Ordinates of Nontrivial Zeros of the Riemann Zeta Function}

\author{R\^otkiv Ayn\"asac}
\email{rotkivanyasac@gmail.com}
\date{June 10, 2026}

\begin{document}

\begin{abstract}
We study exponential sums formed from the ordinates of the nontrivial
zeros of the Riemann zeta function through the geometry of their
centered phase walk. Landau's fixed-real formula at \(X=e\) gives a
logarithmic confinement baseline for
the complete-zero walk
\(Z_N=\sum_{k\le N}e^{\rho_k-\frac12}\), and under the Riemann Hypothesis
(RH) this specializes to the unit phase walk
\(S_N=\sum_{k\le N}e^{i\gamma_k}\). The main structural point is that
centering forces an exact Brownian-bridge covariance kernel,
equivalently the Green kernel of the finite Dirichlet Laplacian. This
turns the centered moment into a weighted pair-correlation statistic and
gives a sine-mode Green decomposition. The same modal language also
expresses the Montgomery--GUE form factor as a constant-mode projection
of the Green-modal vector. 
\end{abstract}

\maketitle

\tableofcontents

\bigskip

\section{Introduction}
\label{sec:introduction}

The nontrivial zeros of the Riemann zeta function are written
\(\rho=\beta+i\gamma\). Their ordinates \(\gamma\) enter both the
Riemann--von Mangoldt counting formula and the fine-scale
pair-correlation problem. The present paper studies these ordinates
through the phase walk obtained from the increments \(e^{i\gamma_k}\).
The basic objects are the endpoint sum \(S_N\), its triangular centroid
\(C_N\), and the centered quadratic moment \(M_0(N)\).

The preliminary sections recall the standard zeta-function background:
the definition and analytic continuation of \(\zeta(s)\), the functional
equation, the critical strip, the Riemann Hypothesis, Weyl
equidistribution, Hlawka's theorem for zeta ordinates, and the
Landau--Gonek distinction between resonant prime-power frequencies and
non-resonant continuous frequencies. These are included to fix notation
and to separate classical analytic estimates from the structural observation
made here.

The core construction begins with the complete-zero auxiliary walk

\begin{equation}
Z_N=\sum_{k=1}^{N}e^{\rho_k-\frac12},
\label{eq:introduction_complete_zero_walk}
\end{equation}

with centroid \(\mathcal C_N\) and centered moment
\(\mathcal M_0(N)\). Landau's fixed-real formula at \(X=e\) supplies the
logarithmic confinement baseline. The new observation is that the centered
phase-walk moment realizes this confinement through an exact
Brownian-bridge/Green-kernel pair statistic, and that the same
Green-modal language gives a natural expression of the Montgomery--GUE
form factor.

Thus the logical chain is

\begin{equation}
\begin{aligned}
\text{centered phase-walk geometry}
&\Longrightarrow
\text{Brownian-bridge/Green kernel}
\\
&\Longrightarrow
\text{weighted pair correlation}
\\
&\Longrightarrow
\text{Green-modal expression of the Montgomery--GUE form factor}.
\end{aligned}
\label{eq:introduction_logical_chain}
\end{equation}

A central point of the paper is the separation between the fixed phase walk and the deformed phase-walk family. The original walk \(S_N=S_N(1)\) is the non-resonant Landau section associated with \(X=e\). The Montgomery--GUE statistics are obtained from the scaled endpoint \(S_N(t)\), with \(t=\tau\log(T/2\pi)\). The Green-modal framework places the fixed-frequency geometry and the Montgomery endpoint form factor in a common coordinate system.

The Montgomery--GUE material below is formulated through exact finite
Green-modal identities, formal equivalences, and the corresponding
form-factor target.

\section*{Notation}
\label{sec:notation}
\addcontentsline{toc}{section}{Notation}

\begin{center}
\begin{tabular}{p{0.30\textwidth}p{0.62\textwidth}}
\(\rho_k=\beta_k+i\gamma_k\)
&
nontrivial zeros with positive ordinate, counted with multiplicity
\\
\(\gamma_k\)
&
ordinate of a zero, measured in radians
\\
\(e^{i\gamma_k}\)
&
unit phase associated with \(\gamma_k\)
\\
\(e^{it\gamma_k}\)
&
\(t\)-deformed ordinate phase
\\
\(B_k,\ \theta_k\)
&
lifted winding number and angular representative of \(\gamma_k\)
\\
\(e^{2\pi i x}\)
&
Weyl/modulo-one phase convention
\\
\(e^{\rho_k-\frac12}\)
&
complete-zero phase; equals \(e^{i\gamma_k}\) under RH
\\
\(S_N,\ C_N,\ M_0(N)\)
&
unit-phase endpoint, triangular centroid, and centered moment
\\
\(Z_N,\ \mathcal C_N,\ \mathcal M_0(N)\)
&
complete-zero endpoint, centroid, and centered moment
\\
\(W_N(j,k)\)
&
centering weight in the original ordinal variables
\\
\(r=N-j+1,\ s=N-k+1\)
&
reversed pair indices for \(W_N(j,k)\)
\\
\(r=N-k+1\)
&
reversed single index for endpoint/modal vectors
\\
\(K_N(r,s)\)
&
reversed Brownian-bridge/Green kernel
\\
\(\phi_m(r),\ \lambda_m\)
&
Dirichlet sine modes and Green eigenvalues
\\
\(A_m(N;t),\ \mathcal A_m(N)\)
&
unit-phase \(t\)-modal coefficient and complete-zero modal coefficient
\\
\(u_m\)
&
projection of the constant vector onto the \(m\)-th sine mode
\\
\(v_r,\ w_m\)
&
centroid weight vector and its sine-mode coefficients
\\
\(T_N,\ T_0,\ L_N\)
&
\(T_N=\gamma_N\); \(T_0\) is a representative local-window height; \(L_N=\log(T_N/2\pi)/(2\pi)\)
\\
\(t\)
&
Fourier parameter; smoothed Montgomery scaling:
\(t=\tau\log(T_N/2\pi)+y\)
\\
\(\mathfrak F_N(\tau),\ F_N(\tau)\)
&
unsmoothed and smoothed Montgomery form factors
\\
\(\xi_{jk}\)
&
unfolded pair variable \(\xi_{jk}=\frac{\log(T_0/2\pi)}{2\pi}(\gamma_j-\gamma_k)\)
\\
\(\Sigma_N(\tau)\)
&
smoothed modal covariance matrix
\\
\(N_0(T),\ \widetilde\gamma_k\)
&
smooth Riemann--von Mangoldt main term and its inverse grid
\\
\(d_k,\ R_m(N;t)\)
&
zero displacement from the RVM grid and the corresponding modal residual
\end{tabular}
\end{center}

\subsection*{Phase conventions}

Throughout the paper, ordinates \(\gamma_k\) are real numbers measured
in radians. The basic unit phase is therefore

\begin{equation}
e^{i\gamma_k}.
\end{equation}

For a real Fourier parameter \(t\), the corresponding \(t\)-deformed
phase is

\begin{equation}
e^{it\gamma_k}.
\end{equation}

When uniform distribution modulo one is discussed, the standard Weyl
notation is used:

\begin{equation}
e^{2\pi i m x_n}.
\end{equation}

Thus, if one writes \(x_n=\alpha\gamma_n\) modulo one, Weyl sums have
the form

\begin{equation}
e^{2\pi i m\alpha\gamma_n}.
\end{equation}

Taking \(\alpha=1/(2\pi)\) and \(m=1\) gives exactly

\begin{equation}
e^{i\gamma_n}.
\end{equation}

The complete-zero phase is a different object:

\begin{equation}
e^{\rho_k-\frac12},
\qquad
\rho_k=\beta_k+i\gamma_k.
\end{equation}

Under RH this specializes to

\begin{equation}
e^{\rho_k-\frac12}=e^{i\gamma_k}.
\end{equation}

Thus the paper uses \(e^{i\gamma_k}\) for radian phases,
\(e^{2\pi i x}\) only for modulo-one Weyl notation,
\(e^{it\gamma_k}\) for Fourier deformation, and
\(e^{\rho_k-\frac12}\) for complete-zero expressions.

\section{The Riemann Zeta Function}
\label{sec:zeta_function}

This section contains standard material.
The classical references are
\cite{Titchmarsh1986},
\cite{Edwards2001},
and
\cite{IwaniecKowalski2004}.

\subsection{Definition}
\label{subsec:zeta_definition}

\begin{definition}
\label{def:zeta_function}

The Riemann zeta function is defined for
$\Re(s)>1$
by the absolutely convergent series

\begin{equation}
\zeta(s)
=
\sum_{n=1}^{\infty}
\frac{1}{n^{s}}.
\label{eq:zeta_dirichlet_series}
\end{equation}

\end{definition}

Euler showed that the same function admits the product expansion

\begin{equation}
\zeta(s)
=
\prod_{p}
\left(
1-p^{-s}
\right)^{-1},
\qquad
\Re(s)>1,
\label{eq:euler_product}
\end{equation}

where the product extends over all prime numbers.
Equation
\eqref{eq:euler_product}
establishes the fundamental connection between
$\zeta(s)$
and the distribution of primes.

\subsection{Analytic Continuation}
\label{subsec:analytic_continuation}

A classical theorem of Riemann states that
$\zeta(s)$
admits a meromorphic continuation to the whole complex plane,
with a unique simple pole at

\begin{equation}
s=1.
\label{eq:zeta_pole}
\end{equation}

See
\cite{Edwards2001}
or
\cite{Titchmarsh1986}.

\subsection{Functional Equation}
\label{subsec:functional_equation}

Define

\begin{equation}
\xi(s)
=
\frac12
s(s-1)
\pi^{-s/2}
\Gamma\!\left(\frac{s}{2}\right)
\zeta(s).
\label{eq:xi_definition}
\end{equation}

The completed zeta function satisfies

\begin{equation}
\xi(s)
=
\xi(1-s).
\label{eq:functional_equation}
\end{equation}

Equivalently,

\begin{equation}
\zeta(s)
=
2^{s}
\pi^{\,s-1}
\sin\!\left(\frac{\pi s}{2}\right)
\Gamma(1-s)
\zeta(1-s).
\label{eq:functional_equation_zeta}
\end{equation}

The symmetry encoded by
\eqref{eq:functional_equation}
plays a central role in the study of the zeros of
$\zeta(s)$.

\subsection{Special Values and Basic Identities}
\label{subsec:special_values}

Among the classical identities are

\begin{equation}
\zeta(2)
=
\frac{\pi^{2}}{6},
\label{eq:zeta2}
\end{equation}

and more generally

\begin{equation}
\zeta(2m)
=
(-1)^{m+1}
\frac{B_{2m}(2\pi)^{2m}}
{2(2m)!},
\qquad
m\ge 1,
\label{eq:zeta_even_values}
\end{equation}

where
$B_{2m}$
denotes the Bernoulli numbers.

These identities are standard and may be found in
\cite{Titchmarsh1986}
or
\cite{Edwards2001}.

\subsection{The Zeros of the Zeta Function}
\label{subsec:zeta_zeros}

The zeros of
$\zeta(s)$
are traditionally divided into two classes.

\begin{definition}
\label{def:trivial_zeros}

The trivial zeros are located at

\begin{equation}
s=-2,-4,-6,\ldots.
\label{eq:trivial_zeros}
\end{equation}

\end{definition}

All remaining zeros are called nontrivial zeros.

\begin{definition}
\label{def:nontrivial_zero}

A nontrivial zero is a complex number

\begin{equation}
\rho
=
\beta+i\gamma
\label{eq:rho_definition}
\end{equation}

such that

\begin{equation}
\zeta(\rho)=0,
\label{eq:rho_zero_condition}
\end{equation}

with

\begin{equation}
0<\beta<1.
\label{eq:critical_strip}
\end{equation}

\end{definition}

The region
\eqref{eq:critical_strip}
is known as the critical strip.

The functional equation implies that the nontrivial zeros occur in
symmetric sets

\begin{equation}
\rho,
\quad
\overline{\rho},
\quad
1-\rho,
\quad
1-\overline{\rho}.
\label{eq:zero_symmetry}
\end{equation}

\begin{definition}
\label{def:zeta_zero_counting_function}

For \(T>0\), let

\begin{equation}
N_\zeta(T)
=
\#\{\rho=\beta+i\gamma:\zeta(\rho)=0,\ 0<\beta<1,\ 0<\gamma\leq T\},
\label{eq:zeta_zero_counting_function}
\end{equation}

where zeros are counted with multiplicity.

\end{definition}

The Riemann--von Mangoldt formula gives the asymptotic estimate

\begin{equation}
N_\zeta(T)
=
\frac{T}{2\pi}
\log\!\left(
\frac{T}{2\pi}
\right)
-
\frac{T}{2\pi}
+
O(\log T).
\label{eq:riemann_von_mangoldt}
\end{equation}

See
\cite{Titchmarsh1986}
or
\cite{Edwards2001}.

\section{The Riemann Hypothesis}
\label{sec:riemann_hypothesis}

This section contains standard material.

\begin{conjecture}[Riemann Hypothesis]
\label{conj:RH}

Every nontrivial zero of the Riemann zeta function satisfies

\begin{equation}
\Re(\rho)
=
\frac12.
\label{eq:RH}
\end{equation}

\end{conjecture}

Assuming
\eqref{eq:RH},
all nontrivial zeros can be written as

\begin{equation}
\rho_n
=
\frac12
+
i\gamma_n.
\label{eq:RH_zeros}
\end{equation}

The sequence of positive ordinates is written as

\begin{equation}
0<\gamma_1
\leq
\gamma_2
\leq
\gamma_3
\leq
\cdots,
\label{eq:gamma_ordering}
\end{equation}

with multiplicities counted. This sequence is the principal object
studied throughout this work.

The Riemann Hypothesis is one of the most important open problems in
mathematics and has deep implications for the distribution of prime
numbers.
Comprehensive discussions may be found in
\cite{Titchmarsh1986},
\cite{Edwards2001},
and
\cite{IwaniecKowalski2004}.

\section{Uniform Distribution of Zeta Ordinates}
\label{sec:uniform_distribution_zeta_ordinates}

The study of exponential sums over the ordinates of the nontrivial
zeros is closely related to the theory of uniform distribution modulo
one.

The fundamental tool is Weyl's criterion.

\begin{theorem}[Weyl Criterion]
\label{thm:weyl_criterion}

Let $(x_n)$ be a sequence of real numbers.

Then $(x_n)$ is uniformly distributed modulo one if and only if

\begin{equation}
\lim_{N\to\infty}
\frac{1}{N}
\sum_{n=1}^{N}
e^{2\pi i m x_n}
=
0
\qquad
(m\neq 0).
\label{eq:weyl_criterion}
\end{equation}

\end{theorem}

A natural question is therefore whether the scaled ordinates
\(\alpha\gamma_n\) are uniformly distributed modulo one. The first
systematic investigation of this problem appears in the work of
Rademacher \cite{Rademacher1956}, who studied the height-truncated sums

\begin{equation}
A_\alpha(T)
=
\sum_{0<\gamma\le T}
e^{2\pi i \alpha \gamma}.
\label{eq:rademacher_sum}
\end{equation}

Subsequent work by Hlawka established the corresponding
equidistribution theorem.

\begin{theorem}[Hlawka, 1975]
\label{thm:hlawka}

For every nonzero real number \(\alpha\), the sequence
\(\alpha\gamma_n\) is uniformly distributed modulo one. Equivalently,

\begin{equation}
\lim_{N\to\infty}
\frac{1}{N}
\sum_{n=1}^{N}
e^{2\pi i m \alpha \gamma_n}
=
0
\qquad
(m\neq 0).
\label{eq:hlawka_weyl}
\end{equation}

\end{theorem}

Taking \(\alpha=1/(2\pi)\) and \(m=1\) in
\eqref{eq:hlawka_weyl} gives

\begin{equation}
\lim_{N\to\infty}
\frac{1}{N}
\sum_{n=1}^{N}
e^{i\gamma_n}
=
0.
\label{eq:SN_normalized_cancellation}
\end{equation}

This is the classical normalized cancellation of the ordinal phase sums.
The consequences for the main objects of the paper are recorded after
those objects are defined.

The exceptional frequencies that survive averaging are described by the
Landau--Gonek theory discussed in Subsection~\ref{sec:prime_power_frequencies}.

\subsection{The Landau--Gonek Formula}
\label{sec:prime_power_frequencies}

\begin{definition}[von Mangoldt Function]
\label{def:von_mangoldt_function}

The von Mangoldt function is defined by

\begin{equation}
\Lambda(\ell)
=
\begin{cases}
\log p,
&
\ell=p^{m},
\quad p \text{ prime},
\quad m\ge 1,
\\[1ex]
0,
&
\text{otherwise}.
\end{cases}
\label{eq:von_mangoldt}
\end{equation}

\end{definition}

The function \(\Lambda(\ell)\) therefore detects prime powers and assigns
to each prime power \(p^m\) the weight \(\log p\).


The arithmetic weights which appear in the Landau--Gonek formula
originate in the logarithmic derivative of the Euler product. For
\(\Re(s)>1\), the identity

\begin{equation}
-\frac{\zeta'(s)}{\zeta(s)}
=
\sum_{\ell=1}^{\infty}
\frac{\Lambda(\ell)}{\ell^s}
\label{eq:log_derivative_von_mangoldt_bridge}
\end{equation}

holds. Equivalently, the Chebyshev function

\begin{equation}
\psi(x)
=
\sum_{\ell\leq x}\Lambda(\ell)
\label{eq:chebyshev_function_von_mangoldt_bridge}
\end{equation}

is the von Mangoldt weighted summatory function naturally coupled to
the zeros through the explicit formula. Thus the same von Mangoldt weights which
occur in the logarithmic derivative also determine the arithmetic
frequencies that appear in Landau--Gonek sums over zero ordinates.

\subsection{Chebyshev function in logarithmic frequency}
\label{subsec:chebyshev_logarithmic_frequency}

Let

\begin{equation}
\psi(x)=\sum_{n\le x}\Lambda(n)
\end{equation}

be the Chebyshev function. In the logarithmic variable

\begin{equation}
u=\log x,
\end{equation}

write

\begin{equation}
\Psi(u)=\psi(e^u).
\end{equation}

The classical explicit formula may be written schematically, away from
prime-power discontinuities and with the usual interpretation of the
zero sum, as

\begin{equation}
\Psi(u)
=
e^u
-
\sum_\rho \frac{e^{\rho u}}{\rho}
-
\log(2\pi)
-
\frac12\log(1-e^{-2u}).
\label{eq:chebyshev_log_frequency_explicit_formula}
\end{equation}

Thus the zero contribution to \(\Psi(u)\) is built from the exponential
modes

\begin{equation}
e^{\rho u}.
\end{equation}

Formally differentiating with respect to \(u\) gives

\begin{equation}
\frac{d}{du}\Psi(u)
=
e^u
-
\sum_\rho e^{\rho u}
-
\frac{e^{-2u}}{1-e^{-2u}},
\label{eq:chebyshev_log_frequency_formal_derivative}
\end{equation}

again in the regularized or distributional sense appropriate to the
explicit formula.

Consequently, the spectral quantity

\begin{equation}
\sum_\rho e^{\rho u}
\end{equation}

is the logarithmic-frequency zero side associated with the Chebyshev
function. The special value \(u=1\) corresponds to the continuous
frequency

\begin{equation}
X=e^u=e.
\end{equation}

At this point one obtains the complete-zero phase

\begin{equation}
e^{\rho-\frac12}
=
e^{-\frac12}e^{\rho}.
\end{equation}

Under RH,

\begin{equation}
\rho=\frac12+i\gamma,
\end{equation}

and hence

\begin{equation}
e^{\rho-\frac12}=e^{i\gamma}.
\end{equation}

Thus the ordinate phase sum

\begin{equation}
\sum_{0<\gamma\le T}e^{i\gamma}
\end{equation}

is, under RH, the \(u=1\) specialization of the
logarithmic-frequency zero side of the explicit formula.

This passage should not be read as an ordinary pointwise Fourier
inversion from \(\psi(x)\) to \(S_N\). The explicit formula is a
regularized identity. Height truncation, smoothing, or distributional
interpretation is needed to make the zero sum precise. The ordinal sum
\(S_N\) is then related to the height-truncated sum through the
height-to-ordinal bridge.

This regularized Chebyshev discussion is separate from the later
application of Landau's fixed-real formula at \(X=e\). The Landau input
used for the logarithmic bound is the height-truncated fixed-\(X\)
estimate in Proposition~\ref{prop:landau_fixed_real_formula}; it does
not use a symmetric truncation, smoothing in height, or a distributional
interpretation at \(u=1\).

In this regularized logarithmic-frequency sense, the height-truncated
phase sums

\begin{equation}
S_x(T)
=
\sum_{0<\gamma\leq T}x^{i\gamma},
\qquad x\in\mathbb N,\ x>1,
\label{eq:SxT_definition}
\end{equation}

may be viewed as logarithmic-frequency probes of the zero side of the
explicit formula. Prime powers are distinguished because they are
precisely the support of \(\Lambda\).

Unlike the metric theory of Hlawka,
the Landau--Gonek formula studies special frequencies associated with
prime powers.

For a fixed integer \(x>1\), the behaviour of \eqref{eq:SxT_definition}
depends strongly on the arithmetic nature of \(x\).

It plays a central role in the explicit formulae relating prime numbers
and nontrivial zeros of the Riemann zeta function
\cite{Landau1911,Titchmarsh1986,IwaniecKowalski2004}.

The classical Landau formula is naturally stated for the zero sum

\begin{equation}
L_x(T)
=
\sum_{0<\gamma\le T}
x^\rho.
\label{eq:landau_zero_sum}
\end{equation}

For a fixed integer \(x>1\), Landau's formula gives

\begin{equation}
L_x(T)
=
-
\frac{T}{2\pi}
\Lambda(x)
+
O(\log T).
\label{eq:landau_zero_sum_asymptotic}
\end{equation}

In particular, when \(x\) is a prime power,

\begin{equation}
x=p^m,
\label{eq:prime_power_condition}
\end{equation}

the von Mangoldt term in \eqref{eq:landau_zero_sum_asymptotic} is
nonzero.

Assuming the Riemann Hypothesis, each nontrivial zero has the form

\begin{equation}
\rho
=
\frac{1}{2}
+
i\gamma.
\label{eq:RH_zero_decomposition}
\end{equation}

Therefore

\begin{equation}
x^\rho
=
\sqrt{x}\,x^{i\gamma}.
\label{eq:xrho_phase_decomposition}
\end{equation}

Thus the passage from Landau's zero sum in \(x^\rho\) to the pure phase
sum in \(x^{i\gamma}\) introduces the factor \(x^{-1/2}\).


The modern quantitative form of this phenomenon is due to Gonek
\cite{Gonek1993},
who refined Landau's original argument using the explicit formula.

\begin{theorem}[Landau--Gonek, phase form under RH]
\label{thm:landau_gonek_formula}

For frequencies associated with prime powers,
the pure phase sum \eqref{eq:SxT_definition} admits, under the Riemann
Hypothesis, a main term proportional to the von Mangoldt function
divided by \(\sqrt{x}\).

In the fixed-\(x\) form,

\begin{equation}
\sum_{0<\gamma\le T}
x^{\,i\gamma}
=
-
\frac{T}{2\pi}
\frac{\Lambda(x)}{\sqrt{x}}
+
O(\log T).
\label{eq:landau_gonek_asymptotic}
\end{equation}

\end{theorem}

Consequently,
prime powers define distinguished resonant frequencies for exponential
sums over the ordinates of the nontrivial zeros.

This phenomenon stands in sharp contrast with the metric
equidistribution results discussed earlier in this section.

\begin{remark}
\label{rem:prime_power_resonances}

Theorem~\ref{thm:landau_gonek_formula} shows that exponential sums over
the ordinates of the zeta zeros possess distinguished resonant
frequencies indexed by prime powers.

From the viewpoint of Weyl equidistribution,
these frequencies are exceptional:
they generate a nonzero main term at the Landau--Gonek scale.

The contrast between the generic behaviour described by
Theorem~\ref{thm:hlawka}
and the resonant behaviour described by
Theorem~\ref{thm:landau_gonek_formula}
is one of the central themes underlying the study of exponential sums
over zeta ordinates.

\end{remark}


\section[The Exponential Sum SN]{The Exponential Sum \texorpdfstring{\(S_N\)}{SN}}
\label{sec:SN_definition}

This section introduces the main objects investigated in this paper.

\subsection{Definition of the Phase Walk}
\label{subsec:phase_walk}

For each ordinate
\(\gamma_n\)
we associate the unit complex phase

\begin{equation}
z_n
=
e^{i\gamma_n}.
\label{eq:phase_definition}
\end{equation}

\begin{definition}
\label{def:SN}

The cumulative exponential sum is defined by

\begin{equation}
S_N
=
\sum_{k=1}^{N}
e^{i\gamma_k}.
\label{eq:SN_definition}
\end{equation}

\end{definition}

Equivalently,

\begin{equation}
S_N
=
\sum_{k=1}^{N}
\cos(\gamma_k)
+
i
\sum_{k=1}^{N}
\sin(\gamma_k).
\label{eq:SN_real_imag}
\end{equation}

The sequence

\begin{equation}
S_1,S_2,\ldots,S_N
\label{eq:SN_trajectory}
\end{equation}

defines a trajectory in the complex plane obtained by successively
adding the unit vectors
\(e^{i\gamma_k}\).

\begin{definition}
\label{def:CN}

The centroid of the phase-walk trajectory is defined by

\begin{equation}
C_N
=
\frac{1}{N}
\sum_{n=1}^{N}
S_n.
\label{eq:CN_definition}
\end{equation}

\end{definition}

Using the definition of \(S_n\), the centroid can also be written as

\begin{equation}
C_N
=
\frac{1}{N}
\sum_{n=1}^{N}
\sum_{k=1}^{n}
e^{i\gamma_k}.
\label{eq:CN_double_sum}
\end{equation}

Interchanging the order of summation gives the weighted form

\begin{equation}
C_N
=
\frac{1}{N}
\sum_{k=1}^{N}
(N-k+1)e^{i\gamma_k}.
\label{eq:CN_weighted}
\end{equation}

Thus \(C_N\) is not the simple average of the phases.
It is the average of the partial sums \(S_n\), or equivalently a
triangularly weighted average of the increments \(e^{i\gamma_k}\).

The geometric and distributional properties of the trajectory
\((S_n)\), together with the behaviour of its centroid \(C_N\),
constitute the central subject of the present paper.

\begin{remark}
\label{rem:equidistribution_not_centroid}

Classical equidistribution gives normalized cancellation of the phase
increments. Its precise consequences for the endpoint \(S_N\) and for the
triangular centroid \(C_N\) are isolated in Section~\ref{sec:non_resonant_triangular_lg_structure}.

\end{remark}

\section{A Logarithmic Quadratic Bound for the Centered Phase Walk}
\label{sec:M0_logarithmic_quadratic_bound}

This section separates the height-truncated statement from the ordinal
statement and keeps track of possible multiplicities at the boundary
height. The natural object before imposing RH is the complete-zero phase
\(e^{\rho-1/2}\). Under RH this becomes the unit phase \(e^{i\gamma}\)
used in the original walk \(S_N\).

\begin{definition}[Fixed-real Landau resonance function]
\label{def:fixed_real_landau_resonance_function}

For \(X>1\), define

\begin{equation}
\Lambda_{\mathbb R}(X)
=
\begin{cases}
\log p,
&
X=p^m,\quad p\ {\rm prime},\quad m\in\mathbb N,\\[4pt]
0,
&
\text{otherwise}.
\end{cases}
\label{eq:fixed_real_landau_resonance_function}
\end{equation}

This is the fixed-real Landau resonance function. It agrees with the von
Mangoldt weight at integral prime powers and is zero away from them. It
is not the ordinary von Mangoldt function on the integers, but the
fixed-real extension appearing in Landau's formula.

\end{definition}

\begin{proposition}[Landau's fixed-real formula]
\label{prop:landau_fixed_real_formula}

Let \(X>1\) be fixed and real, and write

\begin{equation}
X^\rho
=
\exp(\rho\log X),
\label{eq:fixed_real_power_convention}
\end{equation}

where \(\log X\) denotes the real logarithm. Then, as \(T\to\infty\),

\begin{equation}
\sum_{0<\gamma\le T}
X^\rho
=
-
\frac{T}{2\pi}
\Lambda_{\mathbb R}(X)
+
O_X(\log T),
\label{eq:landau_fixed_real_formula}
\end{equation}

where the sum is over the nontrivial zeros \(\rho=\beta+i\gamma\) of
\(\zeta(s)\) with positive ordinate, counted with multiplicity. The
implied constant may depend on the fixed value of \(X\). No uniformity in
\(X\) is used in this paper.

Equivalently, for each fixed \(X>1\) there are constants \(C_L(X)>0\)
and \(T_L(X)\ge2\) such that, for all \(T\ge T_L(X)\),

\begin{equation}
\left|
\sum_{0<\gamma\le T}
X^\rho
+
\frac{T}{2\pi}
\Lambda_{\mathbb R}(X)
\right|
\le
C_L(X)\log T.
\label{eq:landau_fixed_real_formula_quantitative}
\end{equation}

\end{proposition}

\begin{proof}
This is Landau's fixed-frequency formula. Landau proved the formula for
fixed \(X>1\). The fixed-real formulation used here is the form recalled
by Kaczorowski, Languasco, and Perelli in the displayed fixed-\(X\)
formula in the introduction of
\cite{KaczorowskiLanguascoPerelli2000}
(J. Kaczorowski, A. Languasco, A. Perelli,
\emph{A note on Landau's formula}, Functiones et Approximatio
Commentarii Mathematici \textbf{28} (2000), 173--186):

\begin{equation}
\sum_{0<\gamma\le T}X^\rho
=
-
\frac{T}{2\pi}
\Lambda_{\mathbb R}(X)
+
O_X(\log T),
\qquad T\to\infty,
\end{equation}

where \(\Lambda_{\mathbb R}(X)=\log p\) if \(X=p^m\), \(p\) prime, and
\(\Lambda_{\mathbb R}(X)=0\) otherwise. Uniform versions in \(X\) belong
to the Landau--Gonek line of results, beginning with Gonek's work. The
present paper uses only the single fixed value \(X=e\), so Landau's
fixed-\(X\) form is sufficient.

The convention \eqref{eq:fixed_real_power_convention} fixes the meaning
of \(X^\rho\) for real \(X>1\). No smoothing, regularization, or
limiting approximation of \(X=e\) by integer frequencies is used in this
application of the fixed-real formula.
\end{proof}

\begin{remark}[No approximation by integer frequencies]
\label{rem:no_integer_approximation_for_e}

The use of \(X=e\) below is not obtained by approximating \(e\) by
integers. Such an approximation would require a uniform Landau--Gonek
theorem in \(X\), because the factor \(X^{i\gamma}\) is sensitive to
\(\log X\) at height \(T\). Instead,
Proposition~\ref{prop:landau_fixed_real_formula} is applied directly at
the single fixed real frequency \(X=e\).

\end{remark}

\begin{remark}[Arithmetic term and error constants]
\label{rem:landau_arithmetic_term_error_constants}

The Landau formula used here is not an error-free identity. It has an
arithmetic main term and an analytic error term. In the notation of
Proposition~\ref{prop:landau_fixed_real_formula},

\begin{equation}
\sum_{0<\gamma\le T}X^\rho
=
-
\frac{T}{2\pi}\Lambda_{\mathbb R}(X)
+
E_L(T;X),
\qquad
|E_L(T;X)|\le C_L(X)\log T
\end{equation}

for \(T\ge T_L(X)\). At the non-resonant fixed value \(X=e\), the
arithmetic term vanishes because \(\Lambda_{\mathbb R}(e)=0\), and the
complete-zero height sum is controlled by the remaining error term:

\begin{equation}
|\mathcal Z(T)|
\le
C_L(e)\log T
\qquad (T\ge T_L(e)).
\end{equation}

Thus \(C_L(e)\) is the fixed-\(X=e\) error constant supplied by the
chosen quantitative Landau formulation; it is not an additional
identity.

\end{remark}

\begin{corollary}[The non-resonant frequency \(X=e\)]
\label{cor:landau_nonresonant_frequency_e}

At the fixed real frequency \(X=e\),

\begin{equation}
\Lambda_{\mathbb R}(e)=0.
\label{eq:lambda_real_e_zero}
\end{equation}

Consequently,

\begin{equation}
\sum_{0<\gamma\le T}
e^{\rho}
=
O(\log T),
\label{eq:sum_e_rho_log_bound}
\end{equation}

and therefore

\begin{equation}
\mathcal Z(T)
:=
\sum_{0<\gamma\le T}
e^{\rho-\frac12}
=
O(\log T).
\label{eq:complete_zero_height_walk_log_bound}
\end{equation}

\end{corollary}

\begin{proof}
The number \(e\) is not an integral prime power. Hence, by
Definition~\ref{def:fixed_real_landau_resonance_function},

\begin{equation}
\Lambda_{\mathbb R}(e)=0.
\end{equation}

Applying Proposition~\ref{prop:landau_fixed_real_formula} with \(X=e\)
gives

\begin{equation}
\sum_{0<\gamma\le T}e^{\rho}
=
O(\log T).
\end{equation}

Finally,

\begin{equation}
e^{\rho-\frac12}
=
e^{-1/2}e^{\rho},
\end{equation}

so

\begin{equation}
\mathcal Z(T)
=
e^{-1/2}
\sum_{0<\gamma\le T}e^{\rho}
=
O(\log T).
\end{equation}

This proves \eqref{eq:complete_zero_height_walk_log_bound}.
\end{proof}

\begin{remark}[Height truncation and multiplicities]
\label{rem:height_truncation_multiplicity_safe}

Corollary~\ref{cor:landau_nonresonant_frequency_e} gives the primary
height-truncated estimate. It counts every zero at a given ordinate with
its multiplicity. The ordinal theorem below passes from height
truncation to \(Z_N\) by isolating the boundary block of zeros at the
last ordinate.

\end{remark}

\begin{definition}[Ordinal complete-zero walk]
\label{def:ordinal_complete_zero_walk}

Let \(\rho_k=\beta_k+i\gamma_k\) denote the nontrivial zeros with
positive ordinate, ordered by nondecreasing ordinate and counted with
multiplicity. Define

\begin{equation}
Z_N
=
\sum_{k=1}^{N}e^{\rho_k-\frac12},
\label{eq:ZN_definition}
\end{equation}

Its centroid is

\begin{equation}
\mathcal C_N
=
\frac1N
\sum_{n=1}^{N}Z_n,
\label{eq:calC_definition}
\end{equation}

and its centered quadratic moment is

\begin{equation}
\mathcal M_0(N)
=
\frac1N
\sum_{n=1}^{N}
\left|Z_n-\mathcal C_N\right|^2.
\label{eq:calM0_definition}
\end{equation}

\end{definition}

\begin{lemma}[Multiplicity at a fixed ordinate]
\label{lem:multiplicity_fixed_ordinate}

Let \(m(\Gamma)\) denote the number of nontrivial zeros with positive
ordinate equal to \(\Gamma\), counted with multiplicity. There is an
absolute constant \(C_{\rm RVM}>0\), coming only from the
Riemann--von Mangoldt estimate in unit intervals, such that for all
\(\Gamma\ge2\),

\begin{equation}
m(\Gamma)\le C_{\rm RVM}\log \Gamma.
\label{eq:fixed_ordinate_multiplicity_bound}
\end{equation}

\end{lemma}

\begin{proof}
It suffices to consider \(\Gamma\) large; bounded values of \(\Gamma\)
are absorbed into the implied constant. By definition,

\begin{equation}
m(\Gamma)
\le
N_\zeta(\Gamma+1)-N_\zeta(\Gamma-1).
\end{equation}

The Riemann--von Mangoldt formula gives, for some absolute
\(C_{\rm RVM}\) and all \(T\ge2\),

\begin{equation}
N_\zeta(T+1)-N_\zeta(T-1)\le C_{\rm RVM}\log T,
\end{equation}

and hence \(m(\Gamma)\le C_{\rm RVM}\log \Gamma\).
\end{proof}

\begin{lemma}[Height-to-ordinal passage with multiplicities, effective form]
\label{prop:height_to_ordinal_with_multiplicities}
\label{thm:ZN_log_bound}

Let the nontrivial zeros with positive ordinate be ordered by
nondecreasing ordinate and counted with multiplicity:

\begin{equation}
\rho_k=\beta_k+i\gamma_k,
\qquad
0<\gamma_1\le\gamma_2\le\cdots.
\end{equation}

Define

\begin{equation}
Z_N
=
\sum_{k=1}^{N}e^{\rho_k-\frac12}.
\end{equation}

There are positive constants
\(C_L(e),C_{\rm RVM},C_{\rm ord},T_0,N_0\). Put

\begin{equation}
C_Z
=
C_{\rm ord}\left(C_L(e)+e^{1/2}C_{\rm RVM}\right),
\label{eq:CZ_constant_definition}
\end{equation}

Then, for every \(N\ge N_0\),

\begin{equation}
|Z_N|\le C_Z\log N.
\label{eq:ZN_log_bound_quantitative}
\end{equation}

In particular,

\begin{equation}
Z_N
=
O(\log N).
\label{eq:ZN_log_bound_result}
\end{equation}

\end{lemma}

\begin{proof}
By Landau's fixed-real formula at \(X=e\), there are constants
\(C_L(e)>0\) and \(T_L\ge2\) such that, for all \(T\ge T_L\),

\begin{equation}
\left|
\mathcal Z(T)
\right|
=
\left|
\sum_{0<\Im\rho\le T}e^{\rho-\frac12}
\right|
\le
C_L(e)\log T.
\label{eq:effective_landau_e_height_bound}
\end{equation}

Let

\begin{equation}
\Gamma_N=\gamma_N.
\end{equation}

For this proof write \(T_N=\Gamma_N\). Let

\begin{equation}
I_N=\{k:\gamma_k=T_N\}
\end{equation}

be the full multiplicity block at height \(T_N\), and let
\(J_N\subseteq I_N\) be the part of that block included among the first
\(N\) zeros. Since the positive ordinates form a discrete multiset, one
may choose \(T_N^-<T_N\) with no zero ordinate in
\((T_N^-,T_N)\). Then the height-truncated sum up to \(T_N^-\)
contains exactly the first zeros before the boundary block, while
\(B_N\) contains exactly the selected part \(J_N\) of the boundary
block:

\begin{equation}
Z_N
=
\mathcal Z(T_N^-)+B_N,
\label{eq:ZN_height_bridge}
\end{equation}

where

\begin{equation}
B_N=\sum_{k\in J_N}e^{\rho_k-\frac12}.
\label{eq:boundary_block_partial_sum}
\end{equation}

If \(T_N\ge T_0:=\max(T_L,2)\), then
\eqref{eq:effective_landau_e_height_bound} gives

\begin{equation}
|\mathcal Z(T_N^-)|
\le
C_L(e)\log T_N.
\label{eq:ZN_log_TN_bound}
\end{equation}

Moreover,

\begin{equation}
|B_N|
\le
\sum_{k\in J_N}
\left|e^{\rho_k-\frac12}\right|
=
\sum_{k\in J_N}
e^{\beta_k-\frac12}.
\end{equation}

Since \(0<\beta_k<1\),

\begin{equation}
e^{\beta_k-\frac12}\le e^{1/2}.
\end{equation}

Thus

\begin{equation}
|B_N|
\le
e^{1/2}m(T_N)
 \le
e^{1/2}
\left(
N_\zeta(T_N+1)-N_\zeta(T_N-1)
\right)
\le
e^{1/2}C_{\rm RVM}\log T_N.
\end{equation}
Here the last estimate is Lemma~\ref{lem:multiplicity_fixed_ordinate},
which uses only the Riemann--von Mangoldt formula in a unit interval and
does not require simple ordinates.
Therefore

\begin{equation}
|Z_N|
\le
\left(C_L(e)+e^{1/2}C_{\rm RVM}\right)\log T_N.
\label{eq:ZN_log_GammaN_bound}
\end{equation}

Finally, the Riemann--von Mangoldt formula
\eqref{eq:riemann_von_mangoldt} implies \(T_N=\gamma_N\ll N\log N\) for
\(N\) large. Hence there is a constant \(C_{\rm ord}>0\) such that

\begin{equation}
\log T_N
\le
C_{\rm ord}\log N
\label{eq:ZN_log_TN_log_N_bound}
\end{equation}

for every \(N\ge N_0\), after increasing \(N_0\) if necessary so that
\(T_N\ge T_0\). Substituting this estimate in
\eqref{eq:ZN_log_GammaN_bound} gives
\eqref{eq:ZN_log_bound_quantitative} with \(C_Z\) as in
\eqref{eq:CZ_constant_definition}.

The \(O(\log N)\) statement follows immediately.
\end{proof}

\begin{remark}[From asymptotic notation to explicit inequalities]
\label{rem:explicit_constants_logarithmic_bound}

The preceding proof is the exact mechanism by which the asymptotic
notation is converted into an explicit inequality. Once constants are
fixed, the conversion is mechanical. Landau's fixed-real formula at
\(X=e\) supplies \(C_L(e)\) and \(T_L\) such that

\begin{equation}
|\mathcal Z(T_N^-)|\le C_L(e)\log T_N
\qquad (T_N\ge T_0).
\end{equation}

The boundary block is controlled by Riemann--von Mangoldt in a unit
interval:

\begin{equation}
|B_N|
\le e^{1/2}m(T_N)
\le e^{1/2}C_{\rm RVM}\log T_N.
\end{equation}

The exact decomposition

\begin{equation}
Z_N=\mathcal Z(T_N^-)+B_N
\end{equation}

therefore gives

\begin{equation}
|Z_N|
\le
\left(C_L(e)+e^{1/2}C_{\rm RVM}\right)\log T_N
\end{equation}

for all sufficiently large \(N\). Finally, the ordinal height relation
\(\log T_N\le C_{\rm ord}\log N\), valid for \(N\ge N_0\) after fixing
the ordinal threshold, gives

\begin{equation}
|Z_N|\le C_Z\log N,
\qquad
C_Z=C_{\rm ord}\left(C_L(e)+e^{1/2}C_{\rm RVM}\right).
\end{equation}

Thus effective claims require the constants \(C_L(e)\), \(C_{\rm RVM}\),
\(C_{\rm ord}\), and the threshold \(N_0\) to be fixed. The big-\(O\)
statement records the same chain without specifying their values. Here
\(C_L(e)\) comes from the fixed-\(X=e\) Landau error term,
\(C_{\rm RVM}\) comes from the Riemann--von Mangoldt unit-interval
multiplicity bound, and \(C_{\rm ord}\) comes from the height-to-ordinal
comparison \(\log T_N\le C_{\rm ord}\log N\).

\end{remark}

\begin{corollary}[Complete-zero logarithmic confinement, effective form]
\label{cor:complete_zero_logarithmic_confinement}
\label{cor:calC_log_bound}
\label{thm:calM0_log_square_bound}

With \(C_Z\) and \(N_0\) as in
Lemma~\ref{prop:height_to_ordinal_with_multiplicities}, for every
\(N\ge N_0\),

\begin{equation}
|\mathcal C_N|\le C_Z\log N,
\label{eq:calC_log_bound_quantitative}
\end{equation}

and

\begin{equation}
\mathcal M_0(N)\le 4C_Z^2(\log N)^2.
\label{eq:calM0_log_square_bound_quantitative}
\end{equation}

The factor \(4\) comes only from the triangle inequality
\(|Z_n-\mathcal C_N|\le |Z_n|+|\mathcal C_N|\), giving
\(|Z_n-\mathcal C_N|\le2C_Z\log N\) before squaring.

In particular,

\begin{equation}
\mathcal C_N
=
O(\log N).
\label{eq:calC_log_bound_result}
\end{equation}

\begin{equation}
\mathcal M_0(N)
=
O\!\left((\log N)^2\right).
\label{eq:calM0_log_square_bound}
\end{equation}

\end{corollary}

\begin{proof}
Lemma~\ref{prop:height_to_ordinal_with_multiplicities} gives
\(|Z_n|\le C_Z\log n\) for all \(n\ge N_0\). Increasing \(N_0\), if
necessary, to absorb the finitely many values \(n<N_0\), we have, for
every \(N\ge N_0\) and every \(1\le n\le N\),

\begin{equation}
|Z_n|\le C_Z\log N.
\label{eq:Zn_uniform_log_bound}
\end{equation}

By definition,

\begin{equation}
|\mathcal C_N|
\leq
\frac1N
\sum_{n=1}^{N}|Z_n|
\le
C_Z\log N.
\label{eq:calC_average_log_bound}
\end{equation}

This proves \eqref{eq:calC_log_bound_quantitative} and hence
\eqref{eq:calC_log_bound_result}. Thus

\begin{equation}
|\mathcal C_N|
\le
C_Z\log N.
\label{eq:calC_uniform_log_bound}
\end{equation}

Together with \eqref{eq:Zn_uniform_log_bound}, this gives

\begin{equation}
|Z_n-\mathcal C_N|
\leq
|Z_n|+|\mathcal C_N|
\le
2C_Z\log N,
\qquad
1\leq n\leq N.
\label{eq:complete_zero_centered_uniform_log_bound}
\end{equation}

Squaring and averaging gives
\eqref{eq:calM0_log_square_bound_quantitative}, and hence
\eqref{eq:calM0_log_square_bound}.
\end{proof}

\subsection{Specialization under the Riemann Hypothesis}
\label{subsec:M0_RH_specialization}

\begin{corollary}[Specialization under RH]
\label{cor:RH_specialization_complete_zero_walk}
\label{cor:M0_log_square_bound_RH}

Assume RH. Then

\begin{equation}
\rho_k
=
\frac12+i\gamma_k,
\label{eq:rho_k_RH_specialization}
\end{equation}

and hence

\begin{equation}
e^{\rho_k-\frac12}
=
e^{i\gamma_k}.
\label{eq:complete_zero_phase_RH_specialization}
\end{equation}

Therefore

\begin{equation}
Z_N=S_N,
\qquad
\mathcal C_N=C_N,
\qquad
\mathcal M_0(N)=M_0(N).
\label{eq:complete_zero_RH_identification}
\end{equation}

Consequently,

\begin{equation}
M_0(N)
=
O\!\left((\log N)^2\right).
\label{eq:M0_log_square_bound_RH}
\end{equation}

\end{corollary}

\begin{proof}
The identities follow directly from
\(e^{\rho_k-\frac12}=e^{i\gamma_k}\) under RH and from the definitions of
the complete-zero and unit-phase centroids and centered moments. The
bound for \(M_0(N)\) follows from
Corollary~\ref{cor:complete_zero_logarithmic_confinement}.
\end{proof}

Here

\begin{equation}
M_0(N)
=
\frac1N
\sum_{n=1}^{N}|S_n-C_N|^2.
\label{eq:M0_definition}
\end{equation}

\subsection{Roadmap of the logarithmic bound}
\label{subsec:logarithmic_bound_roadmap}

The logic of this section is as follows. Landau's fixed-real formula is
applied directly at the single non-resonant value \(X=e\), giving the
height-truncated estimate for \(\mathcal Z(T)\). The
height-to-ordinal Lemma~\ref{prop:height_to_ordinal_with_multiplicities}
then passes from \(T\)-truncation to \(N\)-truncation by isolating the
boundary block \(B_N\) and bounding its multiplicity through
Riemann--von Mangoldt. This gives \(Z_N=O(\log N)\) with constants
coming only from Landau at \(X=e\) and from the unit-interval RVM bound.
Averaging the prefix sums yields the complete-zero centroid and centered
moment estimates. Under RH these complete-zero estimates specialize to
the unit-phase objects \(S_N\), \(C_N\), and \(M_0(N)\). The remaining
subsections show that the centered moment is exactly a
Brownian-bridge/Green-kernel weighted pair statistic.

\subsection{Pair-correlation form}
\label{subsec:M0_pair_correlation_form}

The preceding estimate has an exact weighted pair-correlation form. Since

\begin{equation}
Z_n
=
\sum_{k=1}^{n}e^{\rho_k-\frac12},
\label{eq:Zn_prefix_for_pair_form}
\end{equation}

we have

\begin{equation}
|Z_n|^2
=
\sum_{j,k\leq n}
e^{\rho_j-\frac12}e^{\overline{\rho_k}-\frac12}.
\label{eq:Zn_square_pair_expansion}
\end{equation}

Moreover,

\begin{equation}
\mathcal M_0(N)
=
\frac1N
\sum_{n=1}^{N}|Z_n|^2
-
|\mathcal C_N|^2.
\label{eq:calM0_uncentered_minus_centroid}
\end{equation}

Using the triangular form of \(\mathcal C_N\),

\begin{equation}
\mathcal C_N
=
\frac1N
\sum_{k=1}^{N}(N-k+1)e^{\rho_k-\frac12},
\label{eq:calC_triangular_form}
\end{equation}

we obtain

\begin{equation}
\mathcal M_0(N)
=
\sum_{j,k=1}^{N}
W_N(j,k)
e^{\rho_j-\frac12}e^{\overline{\rho_k}-\frac12},
\label{eq:calM0_pair_kernel_identity}
\end{equation}

where

\begin{equation}
W_N(j,k)
=
\frac{N-\max(j,k)+1}{N}
-
\frac{(N-j+1)(N-k+1)}{N^2}.
\label{eq:WN_definition}
\end{equation}

Equivalently, setting

\begin{equation}
a=N-j+1,
\qquad
b=N-k+1,
\label{eq:reverse_ordinal_variables_for_WN}
\end{equation}

we may write

\begin{equation}
W_N(j,k)
=
\frac{\min(a,b)}{N}
-
\frac{ab}{N^2}.
\label{eq:WN_brownian_bridge_form}
\end{equation}

The trace of the centering kernel is explicit:

\begin{equation}
\sum_{j=1}^{N}W_N(j,j)
=
\frac{N^2-1}{6N}.
\label{eq:WN_diagonal_trace}
\end{equation}

Combining \eqref{eq:calM0_pair_kernel_identity} with
Corollary~\ref{thm:calM0_log_square_bound}, we obtain

\begin{equation}
\sum_{j,k=1}^{N}
W_N(j,k)e^{\rho_j-\frac12}e^{\overline{\rho_k}-\frac12}
=
O\!\left((\log N)^2\right).
\label{eq:weighted_pair_correlation_bound}
\end{equation}

This is a Brownian-bridge weighted pair-correlation cancellation for the
complete-zero walk.

Here and below, the phrase ``Brownian-bridge'' refers to the exact
algebraic equality of the centering kernel with the finite
Brownian-bridge covariance kernel. No probabilistic approximation, no
Gaussian process limit, and no distributional convergence statement is
being asserted at this point.

Under RH, the pair-correlation identity
\eqref{eq:calM0_pair_kernel_identity} becomes

\begin{equation}
M_0(N)
=
\sum_{j,k=1}^{N}
W_N(j,k)e^{i(\gamma_j-\gamma_k)}.
\label{eq:M0_pair_kernel_identity_RH}
\end{equation}

In particular,

\begin{equation}
\sum_{j,k=1}^{N}
W_N(j,k)e^{i(\gamma_j-\gamma_k)}
=
O\!\left((\log N)^2\right).
\label{eq:weighted_pair_correlation_bound_RH}
\end{equation}

\begin{lemma}[Diagonal and off-diagonal balance]
\label{lem:diagonal_off_diagonal_balance}

Assume RH, so that the phase increments have unit modulus. Then

\begin{equation}
M_0(N)=D_N+O_N,
\end{equation}

where

\begin{equation}
D_N
=
\sum_{j=1}^{N}W_N(j,j)
=
\frac{N^2-1}{6N},
\label{eq:diagonal_contribution_DN_RH}
\end{equation}

and

\begin{equation}
O_N
=
\sum_{\substack{1\le j,k\le N\\j\ne k}}
W_N(j,k)e^{i(\gamma_j-\gamma_k)}.
\label{eq:off_diagonal_contribution_ON_RH}
\end{equation}

Consequently, if
\(M_0(N)=O((\log N)^2)\), then

\begin{equation}
O_N
=
-
\frac{N^2-1}{6N}
+
O\!\left((\log N)^2\right).
\label{eq:off_diagonal_cancellation_law_RH}
\end{equation}

\end{lemma}

\begin{proof}
Under RH,

\begin{equation}
|e^{i\gamma_j}|=1,
\end{equation}

and hence the diagonal part of the pair statistic is

\begin{equation}
D_N
=
\sum_{j=1}^{N}W_N(j,j).
\end{equation}

From

\begin{equation}
W_N(j,k)
=
\frac{N-\max(j,k)+1}{N}
-
\frac{(N-j+1)(N-k+1)}{N^2},
\end{equation}

we get

\begin{equation}
W_N(j,j)
=
\frac{N-j+1}{N}
-
\frac{(N-j+1)^2}{N^2}.
\end{equation}

Put

\begin{equation}
a=N-j+1.
\end{equation}

As \(j\) runs from \(1\) to \(N\), \(a\) runs from \(N\) down to \(1\).
Therefore

\begin{equation}
D_N
=
\sum_{a=1}^{N}
\left(
\frac{a}{N}
-
\frac{a^2}{N^2}
\right).
\end{equation}

Using

\begin{equation}
\sum_{a=1}^{N}a=\frac{N(N+1)}2,
\qquad
\sum_{a=1}^{N}a^2=\frac{N(N+1)(2N+1)}6,
\end{equation}

we obtain

\begin{equation}
D_N
=
\frac{N+1}{2}
-
\frac{(N+1)(2N+1)}{6N}
=
\frac{N^2-1}{6N}.
\end{equation}

Since \(M_0(N)=D_N+O_N\), the claimed off-diagonal identity follows
immediately from \(M_0(N)=O((\log N)^2)\).
\end{proof}

Thus the logarithmic bound on the centered moment is much stronger than
the diagonal size. The diagonal contribution has order \(N\), while the
full centered moment is only \(O((\log N)^2)\). Therefore the
off-diagonal contribution must cancel the diagonal to accuracy
\(O((\log N)^2)\) for this specific Brownian-bridge kernel.

In the complete-zero setting the diagonal contains the factor

\begin{equation}
e^{\rho_j-\frac12}e^{\overline{\rho_j}-\frac12}
=
e^{2\beta_j-1}.
\end{equation}

Thus the complete-zero diagonal is

\begin{equation}
\sum_{j=1}^{N}W_N(j,j)e^{2\beta_j-1}.
\label{eq:complete_zero_diagonal_contribution}
\end{equation}

Only under RH does this reduce to

\begin{equation}
\sum_{j=1}^{N}W_N(j,j)
=
\frac{N^2-1}{6N}.
\end{equation}

Thus the RH specialization recovers the original centered phase walk and
identifies its logarithmic quadratic confinement with Brownian-bridge
weighted off-diagonal cancellation.

\subsection{The Brownian-bridge Green kernel forced by centering}
\label{subsec:what_is_new_present_formulation}

The logarithmic bound follows from Landau's fixed-real formula at \(X=e\),
together with the elementary passage from endpoint bounds to
centered quadratic bounds.

The estimate itself is classical in origin. The structural point of the
present formulation is different: the centered phase-walk moment forces
an exact Green-kernel pair statistic. Thus the centroid geometry
converts a one-dimensional phase walk into a weighted pair-correlation
statistic with Brownian-bridge covariance kernel.

Starting from the complete-zero phase-walk centroid \(\mathcal C_N\),
the centered quadratic moment \(\mathcal M_0(N)\) produces an exact
pair-correlation identity with the explicit weight
\eqref{eq:WN_definition}. After the change of variables
\eqref{eq:reverse_ordinal_variables_for_WN}, this weight becomes
\eqref{eq:WN_brownian_bridge_form}. This is the standard discrete
Brownian-bridge covariance kernel, equivalently the Green kernel of the
discrete Dirichlet Laplacian on the corresponding finite interval.

Thus the structural point is the way the centered geometry of the
zeta-zero phase walk forces exactly this Green/Brownian-bridge kernel and
therefore converts the centroid problem into a precise weighted
pair-correlation statistic of the zeta zeros.

This is an algebraic identity between finite quadratic forms and kernels.
It is not a claim that the zeta-zero phase walk converges to a Brownian
bridge, nor that the modal coefficients are Gaussian. Any probabilistic
or distributional strengthening would require additional hypotheses and
error estimates not used in the present argument.

In short:

\begin{equation}
\begin{aligned}
\text{centered phase-walk geometry}
&\Longrightarrow
\text{Brownian-bridge Green kernel}
\\
&\Longrightarrow
\text{weighted pair correlation of zeta zeros}.
\end{aligned}
\label{eq:phase_walk_green_kernel_pair_correlation}
\end{equation}

Landau's fixed-real formula at \(X=e\) supplies the logarithmic
confinement baseline. The new observation is that the centered
phase-walk moment realizes this
confinement through an exact Brownian-bridge/Green-kernel pair statistic,
and that the same Green-modal language gives a natural expression of the
Montgomery--GUE form factor.

In the RH/unit-phase specialization, the diagonal part of the
Brownian-bridge pair statistic has size \eqref{eq:WN_diagonal_trace}.
The logarithmic bound therefore expresses a near-total off-diagonal
cancellation against this specific Green kernel. The structural
contribution is the identification of the precise kernel and geometric
object through which Landau non-resonance appears in the phase-walk
formulation.

\subsection{Spectral lifting and the Green-kernel modal decomposition}
\label{subsec:spectral_lifting_green_kernel_modal_decomposition}

The centered phase-walk moment is not only a scalar statistic. It is an
exact Hermitian Green-kernel energy, and the Green kernel admits a
complete spectral decomposition.

\begin{lemma}[Lifted phase representation]
\label{lem:lifted_phase_representation}

For every ordinate \(\gamma_k>0\), define

\begin{equation}
B_k=
\left\lfloor \frac{\gamma_k}{2\pi}\right\rfloor,
\qquad
\theta_k=\gamma_k-2\pi B_k\in[0,2\pi).
\label{eq:angular_lifting_ordinate}
\end{equation}

Then

\begin{equation}
\gamma_k=2\pi B_k+\theta_k,
\label{eq:lifting_recovers_ordinate}
\end{equation}

and

\begin{equation}
e^{i\gamma_k}=e^{i\theta_k}.
\label{eq:phase_depends_on_theta}
\end{equation}

Moreover, for every pair \(j,k\),

\begin{equation}
\gamma_j-\gamma_k
=
2\pi(B_j-B_k)+(\theta_j-\theta_k).
\label{eq:lifted_pair_difference}
\end{equation}

Consequently, for every real \(t\),

\begin{equation}
e^{it(\gamma_j-\gamma_k)}
=
e^{i2\pi t(B_j-B_k)}
e^{it(\theta_j-\theta_k)}.
\label{eq:lifted_pair_difference_t_phase}
\end{equation}

\end{lemma}

\begin{proof}
The identity \(\gamma_k=2\pi B_k+\theta_k\) is the Euclidean
decomposition of the real number \(\gamma_k\) modulo \(2\pi\). Since
\(B_k\in\mathbb Z\),

\begin{equation}
e^{i\gamma_k}
=
e^{i(2\pi B_k+\theta_k)}
=
e^{i2\pi B_k}e^{i\theta_k}
=
e^{i\theta_k}.
\end{equation}

Subtracting the decompositions for \(\gamma_j\) and \(\gamma_k\) gives
\eqref{eq:lifted_pair_difference}. Multiplying by \(it\) and
exponentiating gives \eqref{eq:lifted_pair_difference_t_phase}.
\end{proof}

Thus the ordinary phase walk

\begin{equation}
S_N=\sum_{k=1}^{N}e^{i\gamma_k}
\end{equation}

can be formed from the angular representatives alone:

\begin{equation}
S_N=\sum_{k=1}^{N}e^{i\theta_k}.
\end{equation}

However, the angular representatives alone do not determine
\(\gamma_k\). The additional winding integers \(B_k\) restore the full
ordinate:

\begin{equation}
\gamma_k=2\pi B_k+\theta_k.
\end{equation}

For \(t=1\), the winding is invisible:

\begin{equation}
e^{i\gamma_k}=e^{i\theta_k}.
\end{equation}

For general real \(t\),

\begin{equation}
e^{it\gamma_k}
=
e^{i2\pi tB_k}e^{it\theta_k},
\end{equation}

so the winding data are needed to reconstruct the true real-frequency
phase. In finite-\(N\) verification, \(B_k\) is not
an additional hypothesis: it is computed directly from the ordinate by
\eqref{eq:angular_lifting_ordinate}. For Montgomery scaling
\(t=\tau\log(T_N/2\pi)+y\), retaining \(B_k\) is essential because the factor
\(e^{i2\pi tB_k}\) is generally nontrivial.

There are two reversed-index conventions in the paper. For pair
variables \(j,k\), we use

\begin{equation}
r=N-j+1,
\qquad
s=N-k+1.
\label{eq:green_kernel_reversed_indices}
\end{equation}

In these variables the Green/Brownian weight is

\begin{equation}
K_N(r,s)
=
\frac{\min(r,s)}{N}
-
\frac{rs}{N^2},
\qquad
1\le r,s\le N-1.
\label{eq:KN_green_kernel_definition}
\end{equation}

Thus \(W_N(j,k)=K_N(r,s)\). For a single endpoint vector indexed by
\(k\), we use

\begin{equation}
r=N-k+1,
\qquad
x_r(t)=e^{it\gamma_{N-r+1}}.
\end{equation}

These are the same reversal convention applied respectively to a pair
of indices and to a single index.

This is, as a finite matrix identity, the covariance kernel of the
discrete Brownian bridge, equivalently the Green kernel of the discrete
Dirichlet Laplacian on the finite interval. The terminology identifies
the kernel; it does not introduce a random process into the proof.


\begin{lemma}[Orthonormality of the Dirichlet sine basis]
\label{lem:dirichlet_sine_orthonormality}

For \(1\le m,r\le N-1\), define
\begin{equation}
\phi_m(r)
=
\sqrt{\frac{2}{N}}\,
\sin\left(\frac{\pi m r}{N}\right).
\label{eq:dirichlet_sine_modes}
\end{equation}

Then
\begin{equation}
\sum_{r=1}^{N-1}
\phi_m(r)\overline{\phi_\ell(r)}
=
\delta_{m\ell}.
\label{eq:dirichlet_sine_orthonormality}
\end{equation}

Since the functions \(\phi_m\) are real-valued, this is equivalently
\[
\sum_{r=1}^{N-1}
\phi_m(r)\phi_\ell(r)
=
\delta_{m\ell}.
\]
Consequently, \(\{\phi_m\}_{m=1}^{N-1}\) is an orthonormal basis of
\(\mathbb C^{N-1}\).

\end{lemma}

\begin{proof}

It is enough to prove
\begin{equation}
\sum_{r=1}^{N-1}
\sin\left(\frac{\pi m r}{N}\right)
\sin\left(\frac{\pi \ell r}{N}\right)
=
\frac{N}{2}\delta_{m\ell}.
\label{eq:unnormalized_sine_orthogonality}
\end{equation}

Using
\[
\sin a\sin b
=
\frac12\bigl(\cos(a-b)-\cos(a+b)\bigr),
\]
we obtain
\begin{equation}
\sum_{r=1}^{N-1}
\sin\left(\frac{\pi m r}{N}\right)
\sin\left(\frac{\pi \ell r}{N}\right)
=
\frac12
\left(
C_{m-\ell}-C_{m+\ell}
\right),
\label{eq:sine_inner_product_Cq_reduction}
\end{equation}
where
\begin{equation}
C_q
=
\sum_{r=1}^{N-1}
\cos\left(\frac{\pi q r}{N}\right).
\label{eq:Cq_cosine_sum_definition}
\end{equation}

We compute \(C_q\) explicitly. Let
\[
\omega_q=e^{i\pi q/N}.
\]
If \(\omega_q\ne1\), then
\begin{equation}
\sum_{r=1}^{N-1}e^{i\pi q r/N}
=
\sum_{r=1}^{N-1}\omega_q^r
=
\frac{\omega_q(1-\omega_q^{N-1})}{1-\omega_q}.
\label{eq:Cq_geometric_sum}
\end{equation}
Taking real parts gives \(C_q\). Equivalently, for \(0\le q\le 2N-1\),
\begin{equation}
C_q
=
\begin{cases}
N-1, & q=0,\\[4pt]
0, & q\ {\rm odd},\\[4pt]
-1, & q\ {\rm even},\ q\ne0,\ 2N.
\end{cases}
\label{eq:Cq_values}
\end{equation}
Also \(C_{-q}=C_q\), since cosine is even.

In the present application,
\[
1\le m,\ell\le N-1.
\]
Hence
\begin{equation}
0\le |m-\ell|\le N-2,
\qquad
2\le m+\ell\le 2N-2.
\label{eq:Cq_index_range}
\end{equation}
Thus the exceptional value \(q=2N\) never occurs in \(m+\ell\), and \(q=0\) occurs only in \(m-\ell\) when \(m=\ell\).

If \(m=\ell\), then \(m-\ell=0\) and \(m+\ell=2m\), with
\[
2\le 2m\le 2N-2.
\]
Therefore
\[
C_{m-\ell}=C_0=N-1,
\qquad
C_{m+\ell}=C_{2m}=-1.
\]
Consequently,
\begin{equation}
\sum_{r=1}^{N-1}
\sin^2\left(\frac{\pi m r}{N}\right)
=
\frac12\bigl((N-1)-(-1)\bigr)
=
\frac{N}{2}.
\label{eq:diagonal_sine_norm}
\end{equation}

If \(m\ne\ell\), there are two cases. If \(m\) and \(\ell\) have different parity, then \(m-\ell\) and \(m+\ell\) are odd, so
\[
C_{m-\ell}=C_{m+\ell}=0.
\]
If \(m\) and \(\ell\) have the same parity, then \(m-\ell\) and \(m+\ell\) are even. Since \(m\ne\ell\), the first is nonzero, and the second is always nonzero. Both lie in the permitted range away from \(2N\), so
\[
C_{m-\ell}=C_{m+\ell}=-1.
\]
In both subcases,
\[
C_{m-\ell}-C_{m+\ell}=0.
\]
Hence
\[
\sum_{r=1}^{N-1}
\sin\left(\frac{\pi m r}{N}\right)
\sin\left(\frac{\pi \ell r}{N}\right)
=
0
\qquad(m\ne\ell).
\]

Multiplying by the normalization factor \(2/N\) gives
\[
\sum_{r=1}^{N-1}
\phi_m(r)\overline{\phi_\ell(r)}
=
\delta_{m\ell},
\]
which is \eqref{eq:dirichlet_sine_orthonormality}. Since there are \(N-1\) orthonormal vectors in the \(N-1\)-dimensional space
\(\mathbb C^{N-1}\), they form an orthonormal basis.

\end{proof}


The Green kernel has the spectral decomposition

\begin{equation}
K_N(r,s)
=
\sum_{m=1}^{N-1}
\lambda_m\phi_m(r)\phi_m(s),
\label{eq:KN_spectral_decomposition}
\end{equation}

where

\begin{equation}
\lambda_m
=
\frac{1}
{4N\sin^2\left(\frac{\pi m}{2N}\right)}.
\label{eq:KN_spectral_eigenvalues}
\end{equation}

\begin{proposition}[Elementary Green-kernel limits]
\label{prop:elementary_green_kernel_limits}

As \(N\to\infty\),

\begin{equation}
\frac{N^2-1}{6N}
\sim
\frac{N}{6}.
\label{eq:WN_trace_asymptotic}
\end{equation}

For each fixed \(m\ge1\),

\begin{equation}
\lambda_m
=
\frac{1}
{4N\sin^2\left(\frac{\pi m}{2N}\right)}
\sim
\frac{N}{\pi^2m^2}.
\label{eq:green_eigenvalue_fixed_m_asymptotic}
\end{equation}

\end{proposition}

\begin{proof}
The first statement is immediate from
\((N^2-1)/(6N)=N/6-1/(6N)\). For fixed \(m\),
\(\sin(\pi m/(2N))\sim \pi m/(2N)\), which gives
\eqref{eq:green_eigenvalue_fixed_m_asymptotic}.
\end{proof}

Under RH, the phase increments are

\begin{equation}
q_k=e^{i\gamma_k}.
\label{eq:RH_phase_increments_qk}
\end{equation}

Define the modal coefficients

\begin{equation}
A_m(N)
=
\sum_{r=1}^{N-1}
\phi_m(r)
q_{N-r+1}.
\label{eq:RH_modal_coefficients_Am}
\end{equation}

Equivalently,

\begin{equation}
A_m(N)
=
\sum_{r=1}^{N-1}
\sqrt{\frac{2}{N}}
\sin\left(\frac{\pi m r}{N}\right)
e^{i\gamma_{N-r+1}}.
\label{eq:RH_modal_coefficients_Am_expanded}
\end{equation}

Then the centered quadratic moment has the exact modal identity

\begin{equation}
M_0(N)
=
\sum_{m=1}^{N-1}
\lambda_m |A_m(N)|^2.
\label{eq:M0_green_kernel_modal_identity}
\end{equation}

For the complete-zero walk, put

\begin{equation}
q_k=e^{\rho_k-\frac12}.
\label{eq:complete_zero_increments_qk}
\end{equation}

Then

\begin{equation}
\mathcal M_0(N)
=
\sum_{m=1}^{N-1}
\lambda_m |\mathcal A_m(N)|^2,
\label{eq:calM0_green_kernel_modal_identity}
\end{equation}

where

\begin{equation}
\mathcal A_m(N)
=
\sum_{r=1}^{N-1}
\phi_m(r)
e^{\rho_{N-r+1}-\frac12}.
\label{eq:complete_zero_modal_coefficients}
\end{equation}

Expanding the squared modulus gives

\begin{equation}
|\mathcal A_m(N)|^2
=
\sum_{r,s=1}^{N-1}
\phi_m(r)\phi_m(s)
e^{\rho_{N-r+1}-\frac12}
e^{\overline{\rho_{N-s+1}}-\frac12}.
\label{eq:complete_zero_modal_energy_expansion}
\end{equation}

Under RH this reduces to the unit phase identity

\begin{equation}
e^{\rho_{N-r+1}-\frac12}
e^{\overline{\rho_{N-s+1}}-\frac12}
=
e^{i\gamma_{N-r+1}}e^{-i\gamma_{N-s+1}}
=
e^{i(\gamma_{N-r+1}-\gamma_{N-s+1})}.
\label{eq:complete_zero_modal_energy_RH_reduction}
\end{equation}

\begin{lemma}[Complete-zero modal Parseval bound]
\label{lem:complete_zero_modal_parseval_bound}

\begin{equation}
\sum_{m=1}^{N-1}|\mathcal A_m(N)|^2
=
\sum_{k=2}^{N}e^{2\beta_k-1}.
\label{eq:complete_zero_modal_parseval_identity}
\end{equation}

\begin{equation}
\sum_{m=1}^{N-1}|\mathcal A_m(N)|^2
\le
e(N-1),
\label{eq:complete_zero_modal_parseval_bound}
\end{equation}

and hence

\begin{equation}
|\mathcal A_m(N)|
\le
\sqrt{e(N-1)}.
\label{eq:complete_zero_modal_bound}
\end{equation}

Under RH this becomes

\begin{equation}
\sum_{m=1}^{N-1}|\mathcal A_m(N)|^2=N-1,
\qquad
|\mathcal A_m(N)|
\le
\sqrt{N-1}.
\label{eq:complete_zero_modal_bound_RH}
\end{equation}

\end{lemma}

\begin{proof}
Apply Parseval to the vector

\begin{equation}
z_r=e^{\rho_{N-r+1}-\frac12}.
\end{equation}

Then

\begin{equation}
\sum_{m=1}^{N-1}|\mathcal A_m(N)|^2
=
\sum_{r=1}^{N-1}
|z_r|^2.
\end{equation}

Since

\begin{equation}
\left|e^{\rho_k-\frac12}\right|^2
=
e^{2\beta_k-1},
\end{equation}

we get

\begin{equation}
\sum_{m=1}^{N-1}|\mathcal A_m(N)|^2
=
\sum_{k=2}^{N}e^{2\beta_k-1}.
\end{equation}

The bound follows from \(0<\beta_k<1\). Under RH, \(\beta_k=1/2\), so
every factor is \(1\).
\end{proof}

This shows that the Brownian-bridge identity is not merely a scalar
rewriting of \(M_0(N)\). It supplies a canonical Hermitian spectral frame
for the centered phase-walk geometry. The scalar moment \(M_0(N)\) is
the total Green energy, while the quantities \(|A_m(N)|^2\) are its
modal energy components.

The lifted representation \eqref{eq:angular_lifting_ordinate} ensures
that the use of the circular phases \(e^{i\gamma_k}\) does not destroy
the real ordinate information, provided the winding integers \(B_k\) are
retained. Therefore the Green-kernel modal representation can be
formulated together with real pair differences
\eqref{eq:lifted_pair_difference}.

\subsection{Green-generated Fourier pair statistics}
\label{subsec:green_generated_fourier_pair_statistics}

The preceding phase walk corresponds to the single frequency \(t=1\).
For pair-correlation analysis one may allow the real Fourier parameter
\(t\). The corresponding Green-weighted statistic is

\begin{equation}
M_0(N;t)
=
\sum_{j,k=1}^{N}
W_N(j,k)e^{it(\gamma_j-\gamma_k)}.
\label{eq:t_deformed_green_statistic}
\end{equation}

For \(t=1\), this recovers the original centered phase-walk statistic.
With

\begin{equation}
A_m(N;t)
=
\sum_{r=1}^{N-1}
\phi_m(r)
e^{it\gamma_{N-r+1}}.
\label{eq:t_modal_coefficient}
\end{equation}

\begin{lemma}[Uniform finite-\(N\) modal bounds]
\label{lem:finite_modal_bounds}

For every \(N\ge2\), every \(1\le m\le N-1\), and every real \(t\),

\begin{equation}
|A_m(N;t)|
\le
\sqrt{N-1}.
\label{eq:unit_phase_modal_bound}
\end{equation}

Moreover,

\begin{equation}
\sum_{m=1}^{N-1}|A_m(N;t)|^2
=
N-1.
\label{eq:unit_phase_modal_parseval_identity}
\end{equation}

Consequently,

\begin{equation}
\frac1N
\sum_{m=1}^{N-1}|A_m(N;t)|^2
\le
1.
\label{eq:unit_phase_modal_normalized_parseval_bound}
\end{equation}

\end{lemma}

\begin{proof}
By Cauchy--Schwarz,

\begin{equation}
|A_m(N;t)|
\le
\left(
\sum_{r=1}^{N-1}|\phi_m(r)|^2
\right)^{1/2}
\left(
\sum_{r=1}^{N-1}|e^{it\gamma_{N-r+1}}|^2
\right)^{1/2}.
\end{equation}

The sine functions \(\phi_m\) form an orthonormal basis of
\(\mathbb C^{N-1}\), hence

\begin{equation}
\sum_{r=1}^{N-1}|\phi_m(r)|^2=1.
\end{equation}

Also,

\begin{equation}
|e^{it\gamma_{N-r+1}}|=1,
\end{equation}

so

\begin{equation}
\sum_{r=1}^{N-1}|e^{it\gamma_{N-r+1}}|^2=N-1.
\end{equation}

This gives \eqref{eq:unit_phase_modal_bound}.

For the second identity, set

\begin{equation}
x_r(t)=e^{it\gamma_{N-r+1}}.
\end{equation}

The coefficients \(A_m(N;t)\) are precisely the sine-basis coefficients
of the vector \(x(t)\). By Parseval,

\begin{equation}
\sum_{m=1}^{N-1}|A_m(N;t)|^2
=
\sum_{r=1}^{N-1}|x_r(t)|^2
=
N-1.
\end{equation}
\end{proof}

Since the estimate is pointwise in \(t\), it applies uniformly to the
Montgomery scaling

\begin{equation}
t=\tau\log(T_N/2\pi)+y.
\end{equation}

Consequently,

\begin{equation}
\int_{-\infty}^{\infty}
e^{-2|y|}
|A_m(N;\tau\log(T_N/2\pi)+y)|^2\,dy
\le
N-1,
\label{eq:modal_weighted_y_integral_bound}
\end{equation}

because

\begin{equation}
\int_{-\infty}^{\infty}e^{-2|y|}\,dy=1.
\end{equation}

The exact modal identity is

\begin{equation}
M_0(N;t)
=
\sum_{m=1}^{N-1}
\lambda_m
|A_m(N;t)|^2.
\label{eq:t_deformed_green_modal_identity}
\end{equation}

This is the \(t\)-deformed Green energy identity. In lifted variables,

\begin{equation}
e^{it(\gamma_j-\gamma_k)}
=
e^{it[2\pi(B_j-B_k)+(\theta_j-\theta_k)]}.
\label{eq:t_family_lifted_fourier_phase}
\end{equation}

Thus the \(t\)-deformed Green statistic is a Fourier statistic of the
lifted real pair differences, not merely of circular phases. It is,
however, Green-weighted. The unweighted Montgomery form factor is
obtained from the endpoint projection identity below, rather than from a
single scalar value of \(M_0(N;t)\).

\subsection{Endpoint projection and Green deweighting}
\label{subsec:endpoint_projection_green_deweighting}

For later Montgomery scaling, retain

\begin{equation}
T_N=\gamma_N,
\qquad
L_N=\frac{\log(T_N/2\pi)}{2\pi}.
\label{eq:form_factor_TN_LN}
\end{equation}

For every real \(t\), define

\begin{equation}
q_k(t)=e^{it\gamma_k}.
\label{eq:form_factor_qk_t}
\end{equation}

In reversed variables,

\begin{equation}
r=N-k+1,
\label{eq:form_factor_reversed_index}
\end{equation}

put

\begin{equation}
x_r(t)
=
q_{N-r+1}(t)
=
e^{it\gamma_{N-r+1}},
\qquad
1\le r\le N-1.
\label{eq:form_factor_reversed_phase_vector}
\end{equation}

The Green/Brownian sine basis is

\begin{equation}
\phi_m(r)
=
\sqrt{\frac2N}
\sin\left(\frac{\pi m r}{N}\right),
\qquad
1\le m,r\le N-1.
\label{eq:form_factor_sine_basis}
\end{equation}

The modal space is \(\mathbb C^{N-1}\): the endpoint projection below
separates the first term \(e^{it\gamma_1}\), and the modal vector is built
from \(x_r(t)=e^{it\gamma_{N-r+1}}\), \(1\le r\le N-1\), corresponding to
the ordinates \(k=2,\ldots,N\).

The modal coefficients are

\begin{equation}
A_m(N;t)
=
\sum_{r=1}^{N-1}
\phi_m(r)x_r(t)
=
\sum_{r=1}^{N-1}
\phi_m(r)e^{it\gamma_{N-r+1}}.
\label{eq:form_factor_modal_coefficients}
\end{equation}

Let

\begin{equation}
A(N;t)
=
(A_1(N;t),\ldots,A_{N-1}(N;t))^T.
\label{eq:green_modal_vector_A}
\end{equation}

The Green-modal energy is

\begin{equation}
\mathcal E_N(t)
=
A(N;t)^*\Lambda_NA(N;t)
=
\sum_{m=1}^{N-1}
\lambda_m|A_m(N;t)|^2,
\label{eq:green_modal_energy_EN}
\end{equation}

where

\begin{equation}
\Lambda_N=\operatorname{diag}(\lambda_1,\ldots,\lambda_{N-1})
\label{eq:green_modal_lambda_matrix}
\end{equation}

and

\begin{equation}
\lambda_m
=
\frac{1}{4N\sin^2\left(\frac{\pi m}{2N}\right)}.
\label{eq:form_factor_green_eigenvalues}
\end{equation}

In the notation of the centered \(t\)-phase walk,

\begin{equation}
\mathcal E_N(t)=M_0(N;t).
\label{eq:green_modal_energy_equals_M0}
\end{equation}

\begin{lemma}[Exact modal inversion]
\label{lem:exact_modal_inversion}

Let

\begin{equation}
x_r(t)=e^{it\gamma_{N-r+1}},
\qquad
1\le r\le N-1,
\end{equation}

and define

\begin{equation}
A_m(N;t)
=
\sum_{r=1}^{N-1}
\phi_m(r)x_r(t).
\end{equation}

Then

\begin{equation}
x_r(t)
=
\sum_{m=1}^{N-1}
A_m(N;t)\phi_m(r),
\qquad
1\le r\le N-1.
\label{eq:exact_modal_inversion}
\end{equation}

\end{lemma}

\begin{proof}
The result is the expansion of \(x(t)\in\mathbb C^{N-1}\) in the
orthonormal basis \(\{\phi_m\}_{m=1}^{N-1}\). Since the \(\phi_m\) are
real-valued,

\begin{equation}
A_m(N;t)
=
\sum_{r=1}^{N-1}\phi_m(r)x_r(t)
\end{equation}

are precisely the basis coefficients. Hence
\eqref{eq:exact_modal_inversion} follows.
\end{proof}

\begin{lemma}[General modal deweighting]
\label{lem:general_modal_deweighting}

Let \(a=(a_1,\ldots,a_{N-1})^T\in\mathbb C^{N-1}\), and define

\begin{equation}
\eta_m
=
\sum_{r=1}^{N-1}\phi_m(r)a_r,
\qquad
1\le m\le N-1.
\label{eq:general_deweighting_eta_definition}
\end{equation}

Then, for every real \(t\),

\begin{equation}
a^*x(t)
=
\eta^*A(N;t).
\label{eq:general_deweighting_exact_recovery}
\end{equation}

Equivalently,

\begin{equation}
a^*x(t)
=
(\Lambda_N^{-1}\eta)^*\Lambda_NA(N;t).
\label{eq:general_deweighting_green_metric_recovery}
\end{equation}

Consequently,

\begin{equation}
|a^*x(t)|^2
\le
(\eta^*\Lambda_N^{-1}\eta)\,\mathcal E_N(t).
\label{eq:general_deweighting_green_bound}
\end{equation}

\end{lemma}

\begin{proof}
By Lemma~\ref{lem:dirichlet_sine_orthonormality}, both \(x(t)\) and
\(a\) have exact sine-basis expansions,

\begin{equation}
x(t)=\sum_{m=1}^{N-1}A_m(N;t)\phi_m,
\qquad
a=\sum_{m=1}^{N-1}\eta_m\phi_m.
\end{equation}

Taking the Hermitian inner product gives
\eqref{eq:general_deweighting_exact_recovery}. Since

\begin{equation}
\eta^*A(N;t)
=
(\Lambda_N^{-1}\eta)^*\Lambda_NA(N;t),
\end{equation}

\eqref{eq:general_deweighting_green_metric_recovery} follows. Applying
Cauchy--Schwarz in the \(\Lambda_N\)-metric gives

\begin{equation}
|\eta^*A(N;t)|^2
\le
(\eta^*\Lambda_N^{-1}\eta)
(A(N;t)^*\Lambda_NA(N;t)),
\end{equation}

which is \eqref{eq:general_deweighting_green_bound}.
\end{proof}

\begin{lemma}[Constant-mode projection]
\label{lem:constant_mode_projection}

Define

\begin{equation}
u_m
=
\sum_{r=1}^{N-1}\phi_m(r).
\label{eq:constant_mode_projection_um}
\end{equation}

Then

\begin{equation}
u_m
=
\begin{cases}
\displaystyle
\sqrt{\dfrac2N}
\cot\left(\dfrac{\pi m}{2N}\right),
&
m\ {\rm odd},
\\[8pt]
0,
&
m\ {\rm even}.
\end{cases}
\label{eq:constant_mode_coefficients_um}
\end{equation}

\end{lemma}

\begin{proof}
Using

\begin{equation}
\sum_{r=1}^{N-1}\sin(rx)
=
\frac{
\sin\left(\frac{(N-1)x}{2}\right)
\sin\left(\frac{Nx}{2}\right)
}{
\sin\left(\frac{x}{2}\right)
},
\label{eq:finite_sine_sum_identity}
\end{equation}

with

\begin{equation}
x=\frac{\pi m}{N},
\end{equation}

we obtain

\begin{equation}
\sum_{r=1}^{N-1}
\sin\left(\frac{\pi mr}{N}\right)
=
\frac{
\sin\left(\frac{\pi m}{2}-\frac{\pi m}{2N}\right)
\sin\left(\frac{\pi m}{2}\right)
}{
\sin\left(\frac{\pi m}{2N}\right)
}.
\end{equation}

If \(m\) is even, then

\begin{equation}
\sin\left(\frac{\pi m}{2}\right)=0,
\end{equation}

and the sum is zero.
If \(m\) is odd, then

\begin{equation}
\sin\left(\frac{\pi m}{2}\right)=(-1)^{(m-1)/2},
\end{equation}

and

\begin{equation}
\sin\left(\frac{\pi m}{2}-\frac{\pi m}{2N}\right)
=
(-1)^{(m-1)/2}
\cos\left(\frac{\pi m}{2N}\right).
\end{equation}

Therefore

\begin{equation}
\sum_{r=1}^{N-1}
\sin\left(\frac{\pi mr}{N}\right)
=
\frac{
\cos\left(\frac{\pi m}{2N}\right)
}{
\sin\left(\frac{\pi m}{2N}\right)
}
=
\cot\left(\frac{\pi m}{2N}\right).
\label{eq:sine_sum_odd_even}
\end{equation}

Multiplying by the normalization factor \(\sqrt{2/N}\) gives
\eqref{eq:constant_mode_coefficients_um}.
\end{proof}

\begin{lemma}[Endpoint recovery from Green-modal coordinates]
\label{cor:endpoint_projection}

Let

\begin{equation}
u_m=\sum_{r=1}^{N-1}\phi_m(r).
\end{equation}

Then

\begin{equation}
\sum_{k=2}^{N}e^{it\gamma_k}
=
\sum_{m=1}^{N-1}
u_mA_m(N;t).
\label{eq:endpoint_without_first_modal_recovery}
\end{equation}

Therefore

\begin{equation}
\sum_{k=1}^{N}e^{it\gamma_k}
=
e^{it\gamma_1}
+
\sum_{m=1}^{N-1}u_mA_m(N;t).
\label{eq:ordinary_exponential_sum_modal_projection}
\end{equation}

Since

\begin{equation}
u_m
=
\begin{cases}
\displaystyle
\sqrt{\dfrac2N}
\cot\left(\dfrac{\pi m}{2N}\right),
&
m\ {\rm odd},
\\[6pt]
0,
&
m\ {\rm even},
\end{cases}
\end{equation}

this may also be written as

\begin{equation}
\sum_{k=1}^{N}e^{it\gamma_k}
=
e^{it\gamma_1}
+
\sum_{\substack{1\le m\le N-1\\m\ {\rm odd}}}
u_mA_m(N;t).
\label{eq:ordinary_exponential_sum_odd_modal_projection}
\end{equation}

\end{lemma}

\begin{proof}
By exact modal inversion,

\begin{equation}
x_r(t)
=
\sum_{m=1}^{N-1}
A_m(N;t)\phi_m(r).
\end{equation}

Summing over \(r\) gives

\begin{equation}
\sum_{r=1}^{N-1}x_r(t)
=
\sum_{m=1}^{N-1}
A_m(N;t)
\sum_{r=1}^{N-1}\phi_m(r)
=
\sum_{m=1}^{N-1}u_mA_m(N;t).
\label{eq:constant_projection_modal_identity}
\end{equation}

Since \(x_r(t)=e^{it\gamma_{N-r+1}}\), the change of variables
\(k=N-r+1\) gives

\begin{equation}
\sum_{r=1}^{N-1}x_r(t)
=
\sum_{k=2}^{N}e^{it\gamma_k}.
\end{equation}

Adding the missing endpoint \(e^{it\gamma_1}\) gives

\begin{equation}
\sum_{k=1}^{N}e^{it\gamma_k}
=
e^{it\gamma_1}
+
\sum_{m=1}^{N-1}u_mA_m(N;t).
\end{equation}

The parity reduction follows from Lemma~\ref{lem:constant_mode_projection}.
\end{proof}

The even modes \(A_m(N;t)\) are not zero in general. They disappear from
the endpoint sum only because their projection coefficient \(u_m\)
vanishes. They still contribute to the centered Green energy

\begin{equation}
M_0(N;t)
=
\sum_{m=1}^{N-1}
\lambda_m|A_m(N;t)|^2.
\label{eq:centered_energy_all_modes_after_endpoint}
\end{equation}

Thus

\begin{equation}
\text{the endpoint projection sees only odd Green modes,}
\label{eq:endpoint_sees_odd_modes}
\end{equation}

whereas

\begin{equation}
\text{the centered Green energy sees all modes.}
\label{eq:centered_energy_sees_all_modes}
\end{equation}

\subsection{Unconditional Green deweighting of the endpoint projection}
\label{subsec:unconditional_green_deweighting_endpoint_projection}
\leavevmode\par\nobreak

Let \(u=(u_1,\ldots,u_{N-1})^T\).

\begin{theorem}[Unconditional finite Green deweighting]
\label{thm:unconditional_finite_green_deweighting}

For every \(N\ge2\) and every real \(t\),

\begin{equation}
\left|
\sum_{k=2}^{N}e^{it\gamma_k}
\right|^2
\le
2N\,\mathcal E_N(t).
\label{eq:finite_green_deweighting}
\end{equation}

Equivalently,

\begin{equation}
\frac1N
\left|
\sum_{k=2}^{N}e^{it\gamma_k}
\right|^2
\le
2\,\mathcal E_N(t).
\label{eq:finite_green_deweighting_normalized}
\end{equation}

\end{theorem}

\begin{proof}
By the endpoint projection identity without the first endpoint,

\begin{equation}
\sum_{k=2}^{N}e^{it\gamma_k}
=
u^*A(N;t).
\label{eq:endpoint_without_first_uA}
\end{equation}

Since \(\Lambda_N\) is positive definite,

\begin{equation}
u^*A(N;t)
=
(\Lambda_N^{-1/2}u)^*
(\Lambda_N^{1/2}A(N;t)).
\end{equation}

By Cauchy--Schwarz,

\begin{equation}
|u^*A(N;t)|^2
\le
(u^*\Lambda_N^{-1}u)
(A(N;t)^*\Lambda_NA(N;t)).
\label{eq:green_deweighting_cauchy}
\end{equation}

It remains to compute \(u^*\Lambda_N^{-1}u\). From
Lemma~\ref{lem:constant_mode_projection},

\begin{equation}
u^*\Lambda_N^{-1}u
=
\sum_{\substack{1\le m\le N-1\\m\ {\rm odd}}}
|u_m|^2\lambda_m^{-1}.
\end{equation}

Since

\begin{equation}
\lambda_m^{-1}
=
4N\sin^2\left(\frac{\pi m}{2N}\right),
\end{equation}

we get, for odd \(m\),

\begin{equation}
|u_m|^2\lambda_m^{-1}
=
\frac2N
\cot^2\left(\frac{\pi m}{2N}\right)
4N
\sin^2\left(\frac{\pi m}{2N}\right)
=
8\cos^2\left(\frac{\pi m}{2N}\right).
\end{equation}

It remains to sum these terms over odd \(m\). Write
\(\cos^2 x=(1+\cos(2x))/2\). For odd \(m\),
\(2x=\pi m/N\). We use the finite cosine sum

\begin{equation}
\sum_{\ell=1}^{J}\cos((2\ell-1)x)
=
\frac{\sin(2Jx)}{2\sin x}.
\end{equation}

If \(N=2J\), then

\begin{equation}
\sum_{\substack{1\le m\le N-1\\m\ {\rm odd}}}
\cos\left(\frac{\pi m}{N}\right)
=
\sum_{\ell=1}^{J}
\cos\left(\frac{(2\ell-1)\pi}{2J}\right)
=
\frac{\sin \pi}{2\sin(\pi/(2J))}
=
0.
\end{equation}

Hence

\begin{equation}
\sum_{\substack{1\le m\le N-1\\m\ {\rm odd}}}
\cos^2\left(\frac{\pi m}{2N}\right)
=
\frac12 J
=
\frac N4.
\end{equation}

If \(N=2J+1\), then

\begin{equation}
\sum_{\substack{1\le m\le N-1\\m\ {\rm odd}}}
\cos\left(\frac{\pi m}{N}\right)
=
\sum_{\ell=1}^{J}
\cos\left(\frac{(2\ell-1)\pi}{2J+1}\right)
=
\frac{\sin(2J\pi/(2J+1))}{2\sin(\pi/(2J+1))}
=
\frac12.
\end{equation}

Therefore

\begin{equation}
\sum_{\substack{1\le m\le N-1\\m\ {\rm odd}}}
\cos^2\left(\frac{\pi m}{2N}\right)
=
\frac12 J+\frac14
=
\frac{2J+1}{4}
=
\frac N4.
\end{equation}

Thus, for every \(N\ge2\),

\begin{equation}
\sum_{\substack{1\le m\le N-1\\m\ {\rm odd}}}
\cos^2\left(\frac{\pi m}{2N}\right)
=
\frac N4,
\end{equation}

and we obtain

\begin{equation}
u^*\Lambda_N^{-1}u=2N.
\end{equation}

Hence

\begin{equation}
\left|
\sum_{k=2}^{N}e^{it\gamma_k}
\right|^2
=
|u^*A(N;t)|^2
\le
2N\,A(N;t)^*\Lambda_NA(N;t)
=
2N\,\mathcal E_N(t).
\end{equation}
\end{proof}

This theorem is unconditional: it uses only the ordinates \(\gamma_k\),
the finite sine basis, and the Green metric.
As a growth estimate it is generally weaker than the logarithmic endpoint
control obtained from Landau's fixed-real formula at \(X=e\); its role is
to record the exact Green deweighting of the endpoint functional.

The coefficient vector \(A(N;t)\) determines the endpoint projection
exactly through \(u^*A(N;t)\). By contrast, the scalar Green energy

\begin{equation}
\mathcal E_N(t)=A(N;t)^*\Lambda_NA(N;t)
\end{equation}

does not determine the endpoint projection by itself, since it discards
modal direction and phase information. It gives a one-sided estimate, not
an exact inversion. The exact Green-metric recovery identity is

\begin{equation}
u^*A(N;t)
=
(\Lambda_N^{-1}u)^*\Lambda_NA(N;t).
\label{eq:green_metric_endpoint_recovery_identity}
\end{equation}

Thus the scalar Green energy controls the endpoint projection from above.
Exact endpoint recovery uses the coefficient vector \(A(N;t)\),
equivalently the weighted modal data \(\Lambda_NA(N;t)\), through the
linear functional in \eqref{eq:green_metric_endpoint_recovery_identity}.

\begin{corollary}[Centroid deweighting]
\label{cor:centroid_green_deweighting}

For real \(t\), define the \(t\)-deformed triangular centroid

\begin{equation}
C_N(t)
=
\frac1N
\sum_{k=1}^{N}
(N-k+1)e^{it\gamma_k}.
\label{eq:t_deformed_centroid_definition}
\end{equation}

Let \(v=(v_1,\ldots,v_{N-1})^T\) be given in reversed variables by

\begin{equation}
v_r=\frac rN,
\qquad
1\le r\le N-1.
\label{eq:centroid_deweighting_vector_v}
\end{equation}

and define its sine coefficients by

\begin{equation}
w_m
=
\sum_{r=1}^{N-1}v_r\phi_m(r).
\label{eq:centroid_deweighting_coefficients_wm}
\end{equation}

Then

\begin{equation}
C_N(t)-e^{it\gamma_1}
=
w^*A(N;t).
\label{eq:centroid_exact_modal_recovery}
\end{equation}

Moreover,

\begin{equation}
w^*\Lambda_N^{-1}w
=
N-1,
\label{eq:centroid_green_dual_norm}
\end{equation}

and therefore

\begin{equation}
|C_N(t)-e^{it\gamma_1}|^2
\le
(N-1)\mathcal E_N(t).
\label{eq:centroid_green_deweighting_bound}
\end{equation}

\end{corollary}

\begin{proof}
Using \(r=N-k+1\) for \(k=2,\ldots,N\), the definition
\eqref{eq:t_deformed_centroid_definition} gives

\begin{equation}
C_N(t)-e^{it\gamma_1}
=
\sum_{r=1}^{N-1}\frac rN x_r(t)
=
v^*x(t).
\label{eq:centroid_reversed_linear_functional}
\end{equation}

By exact modal inversion, \(x(t)=\sum_m A_m(N;t)\phi_m\), while
\(v=\sum_m w_m\phi_m\). Hence \(v^*x(t)=w^*A(N;t)\), proving
\eqref{eq:centroid_exact_modal_recovery}.

It remains to compute the Green dual norm. In physical variables,
\eqref{eq:physical_green_inverse} gives \(K_N^{-1}=NL_N\), where \(L_N\)
is the finite Dirichlet Laplacian. Hence

\begin{equation}
w^*\Lambda_N^{-1}w
=
v^*K_N^{-1}v
=
N\,v^*L_Nv.
\label{eq:centroid_green_dual_norm_physical}
\end{equation}

For \(v_r=r/N\),

\begin{equation}
v^*L_Nv
=
v_1^2
+
\sum_{r=1}^{N-2}(v_{r+1}-v_r)^2
+
v_{N-1}^2
=
\frac{N-1}{N}.
\label{eq:centroid_laplacian_quadratic_form}
\end{equation}

Thus \(w^*\Lambda_N^{-1}w=N-1\). Applying Cauchy--Schwarz in the
\(\Lambda_N\)-metric gives

\begin{equation}
|w^*A(N;t)|^2
\le
(w^*\Lambda_N^{-1}w)
(A(N;t)^*\Lambda_NA(N;t))
=
(N-1)\mathcal E_N(t),
\label{eq:centroid_deweighting_cauchy}
\end{equation}

which is \eqref{eq:centroid_green_deweighting_bound}.
\end{proof}

This Green bound is generally weaker than the logarithmic centroid bound
obtained by averaging Landau's endpoint estimate. Its role is structural:
it gives the exact Green deweighting of the triangular centroid
functional.

\begin{remark}[Physical-space Green inverse]
\label{rem:physical_space_green_inverse}

In reversed variables,

\begin{equation}
K_N(r,s)
=
\frac{\min(r,s)}{N}
-
\frac{rs}{N^2},
\qquad
1\le r,s\le N-1.
\end{equation}

This matrix is positive definite and satisfies

\begin{equation}
K_N^{-1}=N L_N,
\label{eq:physical_green_inverse}
\end{equation}

where \(L_N\) is the discrete Dirichlet Laplacian

\begin{equation}
L_N=
\begin{pmatrix}
2&-1&0&\cdots&0\\
-1&2&-1&\cdots&0\\
0&-1&2&\cdots&0\\
\vdots&\vdots&\vdots&\ddots&-1\\
0&0&0&-1&2
\end{pmatrix}.
\end{equation}

For the constant vector \(\mathbf 1=(1,\ldots,1)^T\),

\begin{equation}
\mathbf 1^*K_N^{-1}\mathbf 1
=
N\,\mathbf 1^*L_N\mathbf 1
=
2N.
\end{equation}

Therefore, for every \(x\in\mathbb C^{N-1}\),

\begin{equation}
|\mathbf 1^*x|^2
\le
(\mathbf 1^*K_N^{-1}\mathbf 1)(x^*K_Nx)
=
2N\,x^*K_Nx.
\end{equation}

Taking \(x_r=e^{it\gamma_{N-r+1}}\) gives

\begin{equation}
\left|
\sum_{k=2}^{N}e^{it\gamma_k}
\right|^2
\le
2N\,\mathcal E_N(t).
\end{equation}

This is the operator inversion behind the deweighting estimate.
\end{remark}

Let

\begin{equation}
t=\tau\log(T_N/2\pi)+y.
\label{eq:deweighting_montgomery_t}
\end{equation}

Define the endpoint form factor without the first endpoint by

\begin{equation}
F_N^{(2)}(\tau)
=
\frac1N
\int_{-\infty}^{\infty}
e^{-2|y|}
\left|
\sum_{k=2}^{N}
e^{i(\tau\log(T_N/2\pi)+y)\gamma_k}
\right|^2dy.
\label{eq:smoothed_endpoint_form_factor_without_first}
\end{equation}

\begin{corollary}[Unconditional smoothed deweighting bound]
\label{cor:unconditional_smoothed_deweighting_bound}

For every \(N\ge2\) and every real \(\tau\),

\begin{equation}
F_N^{(2)}(\tau)
\le
2
\int_{-\infty}^{\infty}
e^{-2|y|}
\mathcal E_N(\tau\log(T_N/2\pi)+y)\,dy.
\label{eq:smoothed_deweighting_bound_without_first}
\end{equation}

\end{corollary}

\begin{proof}
Apply Theorem~\ref{thm:unconditional_finite_green_deweighting}
pointwise with \(t=\tau\log(T_N/2\pi)+y\). Then integrate against
\(e^{-2|y|}dy\) and divide by \(N\).
\end{proof}

Define the full smoothed endpoint form factor by

\begin{equation}
F_N(\tau)
=
\frac1N
\int_{-\infty}^{\infty}
e^{-2|y|}
\left|
\sum_{k=1}^{N}
e^{i(\tau\log(T_N/2\pi)+y)\gamma_k}
\right|^2dy.
\label{eq:smoothed_form_factor_integral}
\end{equation}

\begin{corollary}[Full endpoint bound]
\label{cor:full_endpoint_bound}

For every \(N\ge2\), every real \(\tau\), and every \(\eta>0\),

\begin{equation}
F_N(\tau)
\le
2(1+\eta)
\int_{-\infty}^{\infty}
e^{-2|y|}
\mathcal E_N(\tau\log(T_N/2\pi)+y)\,dy
+
\frac{1+\eta^{-1}}{N}.
\label{eq:full_endpoint_deweighting_bound}
\end{equation}

\end{corollary}

\begin{proof}
Use
\(|a+b|^2\le (1+\eta)|a|^2+(1+\eta^{-1})|b|^2\) with

\begin{equation}
t=\tau\log(T_N/2\pi)+y,
\qquad
a=\sum_{k=2}^{N}e^{it\gamma_k},
\qquad
b=e^{it\gamma_1}.
\end{equation}

Here \(|b|=1\), and the \(a\)-term is controlled by
Corollary~\ref{cor:unconditional_smoothed_deweighting_bound}.
\end{proof}

\begin{corollary}[Asymptotic deweighting consequence]
\label{cor:asymptotic_deweighting_consequence}

For every fixed \(\tau\in\mathbb R\),

\begin{equation}
\limsup_{N\to\infty}F_N(\tau)
\le
2
\limsup_{N\to\infty}
\int_{-\infty}^{\infty}
e^{-2|y|}
\mathcal E_N(\tau\log(T_N/2\pi)+y)\,dy,
\label{eq:asymptotic_deweighting_consequence}
\end{equation}

with values allowed in \([0,\infty]\). More generally, if

\begin{equation}
\limsup_{N\to\infty}
\int_{-\infty}^{\infty}
e^{-2|y|}
\mathcal E_N(\tau\log(T_N/2\pi)+y)\,dy
\le L(\tau),
\end{equation}

then

\begin{equation}
\limsup_{N\to\infty}F_N(\tau)
\le
2L(\tau).
\end{equation}

If the preceding bound holds uniformly for \(\tau\) in a compact
interval, then the same \(\limsup\) estimate holds uniformly on that
interval.
\end{corollary}

\begin{proof}
From Corollary~\ref{cor:full_endpoint_bound},

\begin{equation}
F_N(\tau)
\le
2(1+\eta)
\int_{-\infty}^{\infty}
e^{-2|y|}
\mathcal E_N(\tau\log(T_N/2\pi)+y)\,dy
+
\frac{1+\eta^{-1}}{N}.
\end{equation}

Taking \(\limsup_{N\to\infty}\) gives

\begin{equation}
\limsup_{N\to\infty}F_N(\tau)
\le
2(1+\eta)
\limsup_{N\to\infty}
\int_{-\infty}^{\infty}
e^{-2|y|}
\mathcal E_N(\tau\log(T_N/2\pi)+y)\,dy.
\end{equation}

Letting \(\eta\downarrow0\) gives the result.
\end{proof}

The deweighting estimate gives one-sided control of the endpoint
projection by the Green energy. Exact recovery of the endpoint
projection is vectorial:

\begin{equation}
u^*A(N;t)
=
(\Lambda_N^{-1}u)^*
\Lambda_NA(N;t).
\label{eq:vectorial_endpoint_recovery}
\end{equation}

Thus the scalar energy

\begin{equation}
A(N;t)^*\Lambda_NA(N;t)
\end{equation}

controls the size of the endpoint projection, while the weighted modal
data

\begin{equation}
\Lambda_NA(N;t)
\end{equation}

contain the phase information needed for exact recovery.

\subsection{The Green-modal form factor and the Montgomery--GUE target}
\label{subsec:green_modal_form_factor_montgomery_gue_target}

We now connect the Green/Brownian modal framework with the Montgomery
pair-correlation form factor \cite{Montgomery1973}.

For the unsmoothed Montgomery scaling put

\begin{equation}
t=2\pi\tau L_N=\tau\log(T_N/2\pi).
\label{eq:montgomery_scaling_t_tau}
\end{equation}

Then

\begin{equation}
e^{it\gamma_k}
=
e^{i\tau\log(T_N/2\pi)\gamma_k}
=
(T_N/2\pi)^{i\tau\gamma_k}.
\label{eq:montgomery_scaling_phase}
\end{equation}

Define the unsmoothed normalized form factor

\begin{equation}
\mathfrak F_N(\tau)
=
\frac1N
\left|
\sum_{k=1}^{N}
e^{i2\pi\tau L_N\gamma_k}
\right|^2.
\label{eq:unsmoothed_normalized_form_factor}
\end{equation}

Using \eqref{eq:ordinary_exponential_sum_odd_modal_projection}, this
becomes

\begin{equation}
\mathfrak F_N(\tau)
=
\frac1N
\left|
e^{i2\pi\tau L_N\gamma_1}
+
\sum_{\substack{1\le m\le N-1\\ m\ \mathrm{odd}}}
\sqrt{\frac2N}
\cot\left(\frac{\pi m}{2N}\right)
A_m(N;2\pi\tau L_N)
\right|^2.
\label{eq:unsmoothed_form_factor_modal_projection}
\end{equation}

Thus the Montgomery form factor is exactly the squared norm of the
constant Green-modal projection.

The same \(F_N(\tau)\) has the standard Montgomery double-sum form. Let

\begin{equation}
w(x)=\frac{4}{4+x^2}.
\label{eq:montgomery_smoothing_weight}
\end{equation}

We use the Fourier identity

\begin{equation}
w(x)
=
\int_{-\infty}^{\infty}
e^{-2|y|}e^{ixy}\,dy.
\label{eq:montgomery_weight_fourier_identity}
\end{equation}

By \eqref{eq:montgomery_weight_fourier_identity}, the endpoint-integral
definition \eqref{eq:smoothed_form_factor_integral} is equivalent to

\begin{equation}
F_N(\tau)
=
\frac1N
\sum_{j,k=1}^{N}
(T_N/2\pi)^{i\tau(\gamma_j-\gamma_k)}
w(\gamma_j-\gamma_k).
\label{eq:smoothed_montgomery_form_factor}
\end{equation}

Using the endpoint projection identity,

\begin{equation}
F_N(\tau)
=
\frac1N
\int_{-\infty}^{\infty}
e^{-2|y|}
\left|
e^{i(\tau\log(T_N/2\pi)+y)\gamma_1}
+
\sum_{\substack{1\le m\le N-1\\ m\ \mathrm{odd}}}
u_mA_m(N;\tau\log(T_N/2\pi)+y)
\right|^2
dy.
\label{eq:smoothed_form_factor_green_modal}
\end{equation}

This is exact for every \(N\). It is the smoothed Montgomery form factor
written entirely in Green-modal coordinates.

\begin{proposition}[Finite-to-limit equivalence]
\label{prop:finite_to_limit_green_modal_form_factor}

For each \(N\), identity \eqref{eq:smoothed_form_factor_green_modal} is
exact. Consequently, the limiting assertion

\begin{equation}
F_N(\tau)\to \min(|\tau|,1)
\end{equation}

is equivalent to convergence of the corresponding Green-modal endpoint
projection

\begin{equation}
e^{it\gamma_1}+u^*A(N;t),
\qquad
t=\tau\log(T_N/2\pi)+y.
\end{equation}

The elementary modal bounds above give uniform finite-\(N\) control. The
finite Green deweighting theorem gives one-sided endpoint control from
the Green energy, as recorded in
Corollary~\ref{cor:asymptotic_deweighting_consequence}. The
Montgomery--GUE target is then an asymptotic statement about the
endpoint projection, or equivalently about the modal expression in
\eqref{eq:smoothed_form_factor_green_modal}.

\end{proposition}


\subsection{From the deformed phase walk to the sine kernel}
\label{subsec:deformed_phase_walk_sine_kernel}

We reserve \(j,k,n,N\) for ordinal indices of zeros, while \(T\),
\(T_N\), and \(T_0\) denote heights. For a prefix of length \(N\), we
write

\begin{equation}
T_N=\gamma_N.
\label{eq:TN_as_gammaN_deformed_walk}
\end{equation}

For a local window, \(T_0\) denotes a representative height of the
window, for instance its central or mean height. Thus the Montgomery
scaling is

\begin{equation}
t=\tau\log(T_N/2\pi)
\label{eq:prefix_montgomery_scaling_log_TN_over_2pi}
\end{equation}

for prefixes and

\begin{equation}
t=\tau\log(T_0/2\pi)
\label{eq:local_window_montgomery_scaling_log_T0_over_2pi}
\end{equation}

for local windows.

The original phase walk is the fixed section of a deformed family:

\begin{equation}
S_N=S_N(1)=\sum_{k=1}^{N}e^{i\gamma_k},
\qquad
S_N(t)=\sum_{k=1}^{N}e^{it\gamma_k}.
\label{eq:fixed_and_deformed_phase_walks}
\end{equation}

The section \(t=1\) corresponds to the fixed real frequency \(X=e\).
Its role in the present paper is the fixed-frequency non-resonant
confinement supplied by Landau's formula, together with the associated
centroid and Green-weighted centered moment. The Montgomery observable
belongs to the full deformed family \(S_N(t)\), in the local scaling
\eqref{eq:prefix_montgomery_scaling_log_TN_over_2pi} or
\eqref{eq:local_window_montgomery_scaling_log_T0_over_2pi}.

Squaring the deformed endpoint gives all ordered pairs:

\begin{equation}
\begin{aligned}
|S_N(t)|^2
&=
\left(\sum_{j\le N}e^{it\gamma_j}\right)
\left(\sum_{k\le N}e^{-it\gamma_k}\right)  \\
&=
\sum_{j,k\le N}e^{it(\gamma_j-\gamma_k)}.
\end{aligned}
\label{eq:deformed_endpoint_pair_expansion}
\end{equation}

Thus the Montgomery form factor is a statistic of pair differences
\(\gamma_j-\gamma_k\), not only of consecutive gaps.

For a local window of length \(W\), centered near height \(T_0\), define
the unfolded pair variable

\begin{equation}
\xi_{jk}
=
\frac{\log(T_0/2\pi)}{2\pi}
(\gamma_j-\gamma_k).
\label{eq:unfolded_pair_variable_xi_jk}
\end{equation}

With \(t=\tau\log(T_0/2\pi)\), one has

\begin{equation}
e^{it(\gamma_j-\gamma_k)}
=
e^{i2\pi\tau\xi_{jk}}.
\label{eq:deformed_phase_as_unfolded_pair_fourier_phase}
\end{equation}

Consequently, for the window endpoint,

\begin{equation}
\frac1W
\left|
S_W(\tau\log(T_0/2\pi))
\right|^2
=
\frac1W
\sum_{j,k\le W}
e^{i2\pi\tau\xi_{jk}}.
\label{eq:window_endpoint_fourier_pair_correlation}
\end{equation}

Thus \(S_W(t)\) gives the Fourier transform of the unfolded pair
correlation attached to the local window. In Fourier variables the
Montgomery--GUE target is the form factor

\begin{equation}
F(\tau)=\min(|\tau|,1),
\label{eq:montgomery_gue_form_factor_target_deformed_walk}
\end{equation}

while in unfolded pair variables it is the sine-kernel pair correlation

\begin{equation}
R_2(\xi)
=
1-
\left(
\frac{\sin \pi\xi}{\pi\xi}
\right)^2.
\label{eq:sine_kernel_pair_correlation_target_deformed_walk}
\end{equation}

These are dual descriptions of the same pair-correlation structure: the
ramp-plateau is the Fourier-side form factor, and the sine kernel is the
side of unfolded differences.

The endpoint reconstruction identity gives the same statement in
Green-modal coordinates:

\begin{equation}
S_N(t)
=
e^{it\gamma_1}
+
\sum_{\substack{1\le m\le N-1\\ m\ {\rm odd}}}
u_m A_m(N;t).
\label{eq:deformed_endpoint_green_modal_reconstruction}
\end{equation}

Hence the smoothed Montgomery form factor is

\begin{equation}
F_N(\tau)
=
\frac1N
\int_{-\infty}^{\infty}
e^{-2|y|}
\left|
e^{i(\tau\log(T_N/2\pi)+y)\gamma_1}
+
\sum_{\substack{1\le m\le N-1\\ m\ {\rm odd}}}
u_m A_m(N;\tau\log(T_N/2\pi)+y)
\right|^2dy.
\label{eq:deformed_endpoint_green_modal_smoothed_form_factor}
\end{equation}

This uses the standard local normalization \(\log(T_N/2\pi)\), consistent
with the Riemann--von Mangoldt mean spacing.

The Green energy and the Montgomery endpoint projection are different
outputs of the same vector \(A(N;t)\). The Green-energy output is

\begin{equation}
A(N;t)\longmapsto A(N;t)^*\Lambda_N A(N;t)
=
\sum_{m=1}^{N-1}\lambda_m |A_m(N;t)|^2,
\label{eq:green_energy_output_from_modal_vector}
\end{equation}

which is the \(t\)-deformed centered Green moment. The endpoint output is

\begin{equation}
A(N;t)\longmapsto
\frac1N
\left|
e^{it\gamma_1}+u^*A(N;t)
\right|^2,
\label{eq:montgomery_endpoint_output_from_modal_vector}
\end{equation}

which is the Montgomery form-factor observable. The Green energy
\(A^*\Lambda A\) and the Montgomery endpoint projection
\(N^{-1}|e^{it\gamma_1}+u^*A|^2\) are distinct observables. They are two
different outputs of the same modal vector \(A(N;t)\).

This framework separates the fixed-frequency confinement of \(S_N(1)\)
from the Montgomery--GUE behaviour encoded by the deformed family
\(S_N(t)\). The Fourier-side form factor and the unfolded pair-side
sine-kernel are two dual descriptions of the same pair-correlation
structure.

\subsection{Finite modal identities and the scaled modal variable}
\label{subsec:finite_modal_identities_scaled_modal_variable}

\begin{lemma}[Finite modal integration]
\label{lem:finite_modal_integration}

For each fixed \(N\), each fixed \(\tau\), and the smoothed Montgomery
parameter \(t_N(\tau,y)=\tau\log(T_N/2\pi)+y\), the \(y\)-integral in the
smoothed form factor may be interchanged with the finite modal sums in
\eqref{eq:smoothed_form_factor_green_modal}. More explicitly,

\begin{equation}
\int_{-\infty}^{\infty}
e^{-2|y|}
\left|
e^{it_N(\tau,y)\gamma_1}
+
\sum_{\substack{1\le m\le N-1\\ m\ {\rm odd}}}
u_mA_m(N;t_N(\tau,y))
\right|^2dy
\end{equation}

is the integral of a finite sum of continuous functions of \(y\), and
the interchange is justified by finite linearity of the integral. No
interchange of an infinite modal sum with the limit \(N\to\infty\) is
used.

\end{lemma}

\begin{proof}
For fixed \(N\), the modal expansion contains at most \(N-1\) terms. For
each \(m\), the coefficient \(A_m(N;t_N(\tau,y))\) is a finite sum of
exponentials in \(y\). Hence it is continuous and bounded by
\eqref{eq:unit_phase_modal_bound}. Since
\(\int_{\mathbb R}e^{-2|y|}\,dy=1\), every term is integrable against the
smoothing weight. The desired interchange is therefore only the
linearity of integration over a finite sum.
\end{proof}

\begin{lemma}[Smoothed modal \(L^2\) control]
\label{lem:smoothed_modal_L2_control}

For every fixed \(N\), \(\tau\), and \(1\le m\le N-1\),

\begin{equation}
\int_{-\infty}^{\infty}
e^{-2|y|}
|A_m(N;t_N(\tau,y))|^2\,dy
\le
N-1.
\label{eq:single_mode_smoothed_L2_control}
\end{equation}

Moreover,

\begin{equation}
\int_{-\infty}^{\infty}
e^{-2|y|}
\sum_{m=1}^{N-1}|A_m(N;t_N(\tau,y))|^2\,dy
=
N-1.
\label{eq:all_modes_smoothed_L2_control}
\end{equation}

\end{lemma}

\begin{proof}
The first bound is \eqref{eq:modal_weighted_y_integral_bound}. For the
second identity, use Parseval
\eqref{eq:unit_phase_modal_parseval_identity} pointwise in
\(t=t_N(\tau,y)\), then integrate against \(e^{-2|y|}\). The exchange of
the finite sum over \(m\) with the \(y\)-integral is justified by
Lemma~\ref{lem:finite_modal_integration}.
\end{proof}

The Green-modal decomposition is a finite-dimensional identity. For each
fixed \(N\),

\begin{equation}
\mathcal E_N(t)
=
A(N;t)^*\Lambda_NA(N;t)
=
\sum_{m=1}^{N-1}\lambda_m|A_m(N;t)|^2
\end{equation}

contains only \(N-1\) terms. Hence no convergence issue is involved in
the identity itself. Lemma~\ref{lem:finite_modal_integration} records the
only exact interchange between \(y\)-integration and modal summation used
below, while Lemma~\ref{lem:smoothed_modal_L2_control} records the
finite-\(N\) \(L^2\) domination used in the smoothed covariance
estimates.

The passage \(N\to\infty\) is a different question. In the Montgomery
scaling

\begin{equation}
t=\tau\log(T_N/2\pi)+y,
\end{equation}

the local zero spacing satisfies heuristically

\begin{equation}
\gamma_{k+1}-\gamma_k\sim \frac{2\pi}{\log(T_N/2\pi)}.
\end{equation}

Thus

\begin{equation}
t(\gamma_{k+1}-\gamma_k)
\sim
2\pi\tau.
\end{equation}

Consequently, as a function of the index \(k\), the vector

\begin{equation}
e^{it\gamma_k}
\end{equation}

has macroscopic oscillation. In the sine basis

\begin{equation}
\phi_m(r)
=
\sqrt{\frac2N}
\sin\left(\frac{\pi mr}{N}\right),
\end{equation}

such oscillations are naturally represented by modal indices of order
\(N\), not by finitely many fixed modes \(m=O(1)\).

Therefore a low-mode truncation \(m\le M\) with fixed \(M\) is not a
natural approximation to the Montgomery endpoint projection. Any genuine
modal limiting theorem in the Montgomery regime should be formulated
using the scaled modal variable

\begin{equation}
\theta=\frac{m}{N}
\end{equation}

or an equivalent frequency-window formulation.

The unconditional limiting statement obtained in the present paper is
instead the finite Green deweighting consequence

\begin{equation}
F_N(\tau)
\le
2(1+\eta)
\int_{-\infty}^{\infty}
e^{-2|y|}
\mathcal E_N(\tau\log(T_N/2\pi)+y)\,dy
+
\frac{1+\eta^{-1}}{N},
\end{equation}

and hence

\begin{equation}
\limsup_{N\to\infty}F_N(\tau)
\le
2
\limsup_{N\to\infty}
\int_{-\infty}^{\infty}
e^{-2|y|}
\mathcal E_N(\tau\log(T_N/2\pi)+y)\,dy.
\end{equation}

This uses only finite-\(N\) identities and finite-\(N\) inequalities; it
does not require an unproved interchange of an infinite modal sum with
the limit \(N\to\infty\).

\subsection{Smoothed modal covariance and conditional form-factor criteria}
\label{subsec:smoothed_modal_covariance_criteria}

The pointwise bound \eqref{eq:unit_phase_modal_bound} and the Parseval
identity \eqref{eq:unit_phase_modal_parseval_identity} are unconditional
finite-dimensional estimates. They are useful for normalization, but
they do not by themselves determine the Montgomery form factor. The
form factor depends on the correlated modal contraction in the
constant-mode direction.

For fixed \(\tau\), put

\begin{equation}
t_N(\tau,y)=\tau\log(T_N/2\pi)+y.
\label{eq:modal_covariance_t_parameter}
\end{equation}

Define the smoothed modal covariance matrix

\begin{equation}
\Sigma_N(\tau)
=
\int_{-\infty}^{\infty}
e^{-2|y|}
A(N;t_N(\tau,y))A(N;t_N(\tau,y))^*
\,dy.
\label{eq:smoothed_modal_covariance_matrix}
\end{equation}

Thus

\begin{equation}
\Sigma_N(\tau)_{m\ell}
=
\int_{-\infty}^{\infty}
e^{-2|y|}
A_m(N;t_N(\tau,y))
\overline{A_\ell(N;t_N(\tau,y))}
\,dy.
\label{eq:smoothed_modal_covariance_entries}
\end{equation}

\begin{proposition}[Exact modal covariance expression for \(F_N\)]
\label{prop:modal_covariance_expression_for_FN}

One has

\begin{align}
F_N(\tau)
&=
\frac1N
+
\frac{2}{N}
\Re
\int_{-\infty}^{\infty}
e^{-2|y|}
e^{-it_N(\tau,y)\gamma_1}
u^*A(N;t_N(\tau,y))
\,dy
\nonumber\\
&\qquad
+
\frac1N
u^*\Sigma_N(\tau)u.
\label{eq:FN_modal_covariance_expansion}
\end{align}

Moreover, if

\begin{equation}
\frac1N u^*\Sigma_N(\tau)u=O(1),
\label{eq:modal_covariance_boundedness_assumption}
\end{equation}

then the middle term in \eqref{eq:FN_modal_covariance_expansion} is
\(o(1)\), and therefore

\begin{equation}
F_N(\tau)
=
\frac1N u^*\Sigma_N(\tau)u+o(1).
\label{eq:FN_as_modal_covariance_contraction}
\end{equation}

\end{proposition}

\begin{proof}
Starting from \eqref{eq:smoothed_form_factor_green_modal} and using
\(t=t_N(\tau,y)\), expand the square. The first term gives

\begin{equation}
\frac1N
\int_{-\infty}^{\infty}e^{-2|y|}\,dy
=
\frac1N.
\end{equation}

The quadratic modal term is exactly

\begin{equation}
\frac1N
u^*
\left(
\int_{-\infty}^{\infty}
e^{-2|y|}
A(N;t_N(\tau,y))A(N;t_N(\tau,y))^*
\,dy
\right)
u
=
\frac1N u^*\Sigma_N(\tau)u.
\end{equation}

This gives \eqref{eq:FN_modal_covariance_expansion}. For the middle
term, Cauchy--Schwarz gives

\begin{align}
&\left|
\int_{-\infty}^{\infty}
e^{-2|y|}
e^{-it_N(\tau,y)\gamma_1}
u^*A(N;t_N(\tau,y))\,dy
\right|^2
\nonumber\\
&\qquad\le
\left(\int_{-\infty}^{\infty}e^{-2|y|}\,dy\right)
\left(
\int_{-\infty}^{\infty}
e^{-2|y|}
|u^*A(N;t_N(\tau,y))|^2\,dy
\right)
\nonumber\\
&\qquad=
u^*\Sigma_N(\tau)u.
\label{eq:modal_covariance_middle_term_bound}
\end{align}

Hence the middle term is bounded by

\begin{equation}
\frac{2}{N}
\left(u^*\Sigma_N(\tau)u\right)^{1/2}.
\end{equation}

Under \eqref{eq:modal_covariance_boundedness_assumption}, this is
\(O(N^{-1/2})=o(1)\). This proves
\eqref{eq:FN_as_modal_covariance_contraction}.
\end{proof}

\begin{conjecture}[Modal covariance form of Montgomery]
\label{conj:modal_covariance_montgomery}

For each fixed \(\tau\in\mathbb R\), the smoothed modal covariance
contraction satisfies

\begin{equation}
\frac1N
u^*\Sigma_N(\tau)u
\longrightarrow
\min(|\tau|,1).
\label{eq:modal_covariance_montgomery_conjecture}
\end{equation}

\end{conjecture}

\begin{corollary}[Modal covariance criterion for the Montgomery form factor]
\label{cor:modal_covariance_montgomery_criterion}

Fix \(\tau\in\mathbb R\). All convergence in this corollary concerns the
smoothed \(y\)-averaged form factor \(F_N(\tau)\) defined in
\eqref{eq:smoothed_form_factor_integral}, pointwise in this fixed value
of \(\tau\). Suppose that

\begin{equation}
\frac1N
u^*\Sigma_N(\tau)u
\longrightarrow
L(\tau)
\label{eq:modal_covariance_convergence_assumption}
\end{equation}

for a finite limit \(L(\tau)\). Then

\begin{equation}
F_N(\tau)
\longrightarrow
L(\tau).
\label{eq:FN_converges_to_modal_covariance_limit}
\end{equation}

In particular, the Montgomery--GUE form-factor target is equivalent, in
this modal language, to
Conjecture~\ref{conj:modal_covariance_montgomery}, namely to the
conditional modal covariance contraction

\begin{equation}
\frac1N
u^*\Sigma_N(\tau)u
\longrightarrow
\min(|\tau|,1).
\label{eq:modal_covariance_montgomery_target}
\end{equation}

\end{corollary}

\begin{proof}
The convergence assumption \eqref{eq:modal_covariance_convergence_assumption}
implies the boundedness condition
\eqref{eq:modal_covariance_boundedness_assumption}. The conclusion follows
from Proposition~\ref{prop:modal_covariance_expression_for_FN}.
\end{proof}

\begin{remark}[Individual modal bounds versus modal covariance]
\label{rem:individual_bounds_vs_modal_covariance}

The estimates \eqref{eq:unit_phase_modal_bound} and
\eqref{eq:unit_phase_modal_parseval_identity} are unconditional. They do
not determine the Montgomery form factor, because \(F_N(\tau)\) depends
on the rank-one modal contraction \(u^*\Sigma_N(\tau)u\), which includes
correlations between different modes. Thus individual modal bounds
provide finite control, while the limiting form factor is governed by
the smoothed modal covariance in the constant-mode direction \(u\).

\end{remark}

The Green energy gives a useful one-sided consequence. Since
\eqref{eq:finite_green_deweighting} gives

\begin{equation}
|u^*A(N;t)|^2
\le
2N\,\mathcal E_N(t),
\end{equation}

one has

\begin{equation}
\frac1N u^*\Sigma_N(\tau)u
=
\frac1N
\int_{-\infty}^{\infty}
e^{-2|y|}
|u^*A(N;t_N(\tau,y))|^2\,dy
\le
2
\int_{-\infty}^{\infty}
e^{-2|y|}
\mathcal E_N(t_N(\tau,y))\,dy.
\label{eq:modal_covariance_controlled_by_green_energy}
\end{equation}

Consequently, a bound on the smoothed Green energy gives a bound on the
modal covariance contraction and hence on \(F_N(\tau)\), but not exact
recovery of the limiting form factor unless the covariance contraction
itself is controlled.

The Montgomery--GUE target in this notation is

\begin{equation}
F_N(\tau)
\longrightarrow
\min(|\tau|,1).
\label{eq:montgomery_form_factor_target}
\end{equation}

Equivalently, the Green-modal target is

\begin{equation}
\begin{aligned}
&
\frac1N
\int_{-\infty}^{\infty}
e^{-2|y|}
\left|
e^{i(\tau\log(T_N/2\pi)+y)\gamma_1}
+
\sum_{\substack{1\le m\le N-1\\ m\ \mathrm{odd}}}
u_mA_m(N;\tau\log(T_N/2\pi)+y)
\right|^2
dy
\\
&\hspace{5cm}
\longrightarrow
\min(|\tau|,1).
\end{aligned}
\label{eq:green_modal_montgomery_target}
\end{equation}

This is the Green-modal form of the Montgomery form-factor conjecture.

\begin{remark}[Fixed frequencies versus Montgomery scaling]
\label{rem:fixed_landau_frequencies_vs_montgomery_scaling}

Landau fixed-frequency formulas concern fixed frequencies. The
non-integer fixed frequency used in the complete-zero walk is covered by
Proposition~\ref{prop:landau_fixed_real_formula}; at \(X=e\) it gives

\begin{equation}
\sum_{0<\gamma\le T}e^{\rho}
=
O(\log T).
\label{eq:fixed_frequency_landau_in_montgomery_remark}
\end{equation}

For integral prime-power frequencies the fixed-frequency Landau--Gonek
main term has negative sign and is supported at the arithmetic
resonances \(X=p^m\). In the fixed-real notation
\(\Lambda_{\mathbb R}(e)=0\), since \(e\) is not an integral prime power.

Under RH, the \(X=e\) formula gives the corresponding fixed-frequency phase
statement

\begin{equation}
\sum_{0<\gamma\le T}e^{i\gamma}
=
O(\log T).
\label{eq:fixed_frequency_phase_form_montgomery_remark}
\end{equation}

Montgomery's form factor belongs to a different regime. In the present
ordinal normalization the frequency is \(N\)-dependent; one takes

\begin{equation}
X_N=(T_N/2\pi)^\tau,
\qquad
t=\tau\log(T_N/2\pi),
\label{eq:moving_montgomery_frequency}
\end{equation}

After squaring and smoothing, the diagonal prime contributions in the
explicit formula accumulate over the growing range

\begin{equation}
n\le (T_N/2\pi)^{|\tau|}.
\label{eq:moving_montgomery_prime_range}
\end{equation}

At the formal explicit-formula level, this is the standard ramp mechanism

\begin{equation}
F_N(\tau)
\sim
\frac1{\log T_N}
\sum_{n\le (T_N/2\pi)^{|\tau|}}
\frac{\Lambda(n)^2}{n}
\sim
|\tau|,
\qquad
0<|\tau|<1.
\label{eq:fixed_vs_moving_ramp_mechanism}
\end{equation}

The fixed frequency \(X=e\) is non-resonant for the phase walk \(S_N\),
whereas the Montgomery scaling \(X_N=(T_N/2\pi)^\tau\) explores a growing
arithmetic range after squaring and smoothing.

\end{remark}

Thus the form-factor target has ramp

\begin{equation}
F_N(\tau)\to |\tau|,
\qquad
0<|\tau|<1.
\label{eq:montgomery_ramp_limit}
\end{equation}

and saturation

\begin{equation}
F_N(\tau)\to 1,
\qquad
|\tau|\ge 1.
\label{eq:montgomery_saturation_limit}
\end{equation}

Together these give

\begin{equation}
F_N(\tau)\to\min(|\tau|,1).
\label{eq:montgomery_ramp_saturation_combined}
\end{equation}

Finally connect this to the sine kernel. With the convention

\begin{equation}
\widehat f(\tau)
=
\int_{\mathbb R}f(u)e^{-2\pi i\tau u}\,du,
\label{eq:fourier_transform_convention}
\end{equation}

one has

\begin{equation}
\widehat{
\left(
\frac{\sin \pi u}{\pi u}
\right)^2
}
=
(1-|\tau|)_+.
\label{eq:sine_square_fourier_transform}
\end{equation}

Since

\begin{equation}
\min(|\tau|,1)-1
=
-(1-|\tau|)_+,
\label{eq:form_factor_minus_one_triangle}
\end{equation}

the corresponding limiting pair-correlation density is

\begin{equation}
R_2(u)
=
1-
\left(\frac{\sin \pi u}{\pi u}\right)^2.
\label{eq:gue_pair_correlation_density}
\end{equation}

The Green/Brownian framework does not merely rewrite the centered moment
\(M_0(N;t)\). It also supplies a canonical modal representation of the
Montgomery form factor. The Green energy is the diagonal quadratic form

\begin{equation}
A(N;t)^*\Lambda_N A(N;t),
\label{eq:green_energy_matrix_form}
\end{equation}

whereas the Montgomery form factor is the rank-one constant-mode
contraction

\begin{equation}
\frac1N
\left|
e^{it\gamma_1}+u^*A(N;t)
\right|^2.
\label{eq:form_factor_rank_one_contraction}
\end{equation}

Thus the same Green-modal vector \(A(N;t)\) carries both the centered
phase-walk energy and the pair-correlation form factor. In this
language, the Montgomery--GUE form-factor target is the limiting energy
\(\min(|\tau|,1)\) of the constant-mode Green projection in the unfolded
scaling

\begin{equation}
t=\tau\log(T_N/2\pi).
\label{eq:form_factor_unfolded_scaling_final}
\end{equation}

\subsection{Riemann--von Mangoldt background and modal residuals}
\label{subsec:rvm_background_modal_residuals}

The Green-modal representation also separates a smooth counting
background from the genuine zero fluctuations. Define the smooth
Riemann--von Mangoldt main term

\begin{equation}
N_0(T)
=
\frac{T}{2\pi}\log\frac{T}{2\pi}
-
\frac{T}{2\pi}
+
\frac78.
\label{eq:rvm_smooth_counting_function}
\end{equation}

We use RVM as shorthand for Riemann--von Mangoldt in the notation below.
The corresponding RVM grid is the deterministic grid obtained by
inverting this smooth counting function.

Let

\begin{equation}
\widetilde\gamma_k=N_0^{-1}(k),
\label{eq:rvm_grid_ordinate}
\end{equation}

and define the residual displacement

\begin{equation}
d_k=\gamma_k-\widetilde\gamma_k.
\label{eq:rvm_residual_displacement}
\end{equation}

The modal coefficient decomposes as

\begin{equation}
A_m(N;t)
=
A_m^{\mathrm{RVM}}(N;t)+R_m(N;t),
\label{eq:rvm_modal_decomposition}
\end{equation}

where

\begin{equation}
A_m^{\mathrm{RVM}}(N;t)
=
\sum_{r=1}^{N-1}
\phi_m(r)
e^{it\widetilde\gamma_{N-r+1}},
\label{eq:rvm_background_mode}
\end{equation}

and

\begin{equation}
R_m(N;t)
=
\sum_{r=1}^{N-1}
\phi_m(r)
e^{it\widetilde\gamma_{N-r+1}}
\left(
e^{itd_{N-r+1}}-1
\right).
\label{eq:rvm_residual_mode}
\end{equation}

\begin{lemma}[Parseval bound for the RVM modal residual]
\label{lem:rvm_residual_parseval_bound}

Let

\begin{equation}
\widetilde\gamma_k=N_0^{-1}(k),
\qquad
d_k=\gamma_k-\widetilde\gamma_k,
\end{equation}

and define \(R_m(N;t)\) by \eqref{eq:rvm_residual_mode}. Then

\begin{equation}
\sum_{m=1}^{N-1}|R_m(N;t)|^2
=
\sum_{k=2}^{N}
|e^{itd_k}-1|^2.
\label{eq:rvm_residual_parseval_bound}
\end{equation}

Consequently,

\begin{equation}
\sum_{m=1}^{N-1}|R_m(N;t)|^2
\le
4(N-1).
\label{eq:rvm_residual_uniform_bound}
\end{equation}

Moreover,

\begin{equation}
\sum_{m=1}^{N-1}|R_m(N;t)|^2
\le
t^2\sum_{k=2}^{N}|d_k|^2.
\label{eq:rvm_residual_small_displacement_bound}
\end{equation}

\end{lemma}

\begin{proof}
Set

\begin{equation}
b_r(t)
=
e^{it\widetilde\gamma_{N-r+1}}
\left(
e^{itd_{N-r+1}}-1
\right),
\qquad
1\le r\le N-1.
\end{equation}

Then

\begin{equation}
R_m(N;t)
=
\sum_{r=1}^{N-1}\phi_m(r)b_r(t).
\end{equation}

Since

\begin{equation}
\phi_m(r)
=
\sqrt{\frac2N}
\sin\left(\frac{\pi mr}{N}\right),
\qquad
1\le m,r\le N-1,
\end{equation}

is an orthonormal sine basis of \(\mathbb C^{N-1}\), Parseval's identity
gives

\begin{equation}
\sum_{m=1}^{N-1}|R_m(N;t)|^2
=
\sum_{r=1}^{N-1}|b_r(t)|^2.
\end{equation}

Because

\begin{equation}
|e^{it\widetilde\gamma_{N-r+1}}|=1,
\end{equation}

we have

\begin{equation}
|b_r(t)|^2
=
\left|
e^{itd_{N-r+1}}-1
\right|^2.
\end{equation}

Changing variables \(k=N-r+1\), this becomes

\begin{equation}
\sum_{m=1}^{N-1}|R_m(N;t)|^2
=
\sum_{k=2}^{N}|e^{itd_k}-1|^2.
\end{equation}

The bound

\begin{equation}
|e^{ix}-1|\le2
\end{equation}

gives

\begin{equation}
\sum_{m=1}^{N-1}|R_m(N;t)|^2
\le
4(N-1).
\end{equation}

The sharper small-displacement estimate follows from

\begin{equation}
|e^{ix}-1|\le |x|,
\end{equation}

which gives

\begin{equation}
|e^{itd_k}-1|^2
\le
t^2|d_k|^2.
\end{equation}

Summing over \(k\) proves the final bound.
\end{proof}

The estimate \(4(N-1)\) is a universal Parseval bound. It is sharp at
this level of information: without further control on the phases
\(td_k\), the individual quantities \(|e^{itd_k}-1|\) may be as large as
\(2\), so no \(o(N)\) improvement follows from orthogonality alone.

For fixed \(t\), the refinement

\begin{equation}
\sum_{m=1}^{N-1}|R_m(N;t)|^2
\le
t^2\sum_{k=2}^{N}|d_k|^2
\end{equation}

can be stronger if the displacements
\(d_k=\gamma_k-\widetilde\gamma_k\) are square-summably small on average.
In Montgomery scaling, however,

\begin{equation}
t=\tau\log(T_N/2\pi)+y,
\end{equation}

and the products \(td_k\) need not be small. Thus the Parseval estimate
gives uniform modal boundedness after normalization,

\begin{equation}
\frac1N\sum_{m=1}^{N-1}|R_m(N;t)|^2\le4,
\label{eq:rvm_residual_normalized_bound}
\end{equation}

and it shows that the RVM residual has controlled total modal energy.

Thus \(A_m^{\mathrm{RVM}}\) is the calculable smooth background, while
\(R_m\) is the modal remainder caused by replacing the synthetic
Riemann--von Mangoldt grid ordinates \(\widetilde\gamma_k\) with the
actual zero ordinates \(\gamma_k\). The RVM grid alone is rigid after
unfolding; it is not a GUE process. The Green-modal form factor
separates this smooth background from the residual fluctuations that
carry the arithmetic and spectral information of the actual zeros.

\section[Supplementary Exact Product Forms for SN and CN]{Supplementary Exact Product Forms for \texorpdfstring{\(S_N\)}{SN} and \texorpdfstring{\(C_N\)}{CN}}
\label{sec:product_forms_SN_CN}

This supplementary section gives exact algebraic reformulations of the
cumulative sum \(S_N\) and of the centroid \(C_N\) in terms of finite
products associated with the phases \(e^{i\gamma_k}\). These identities
are complementary algebraic structure; the main structural result of the
paper remains the Brownian-bridge/Green-kernel formulation developed
above.

\subsection{A Finite Product Identity}
\label{subsec:finite_product_identity}

We first record the elementary identity on which the product
formulation is based.

\begin{lemma}[Logarithmic coefficients of a finite product]
\label{lem:finite_product_logarithmic_coefficients}

Let \(a_1,\ldots,a_N\) be complex numbers, and define

\begin{equation}
P(z)
=
\prod_{k=1}^{N}
(1-a_k z).
\label{eq:finite_product_definition}
\end{equation}

For \(|z|\) sufficiently small,

\begin{equation}
\log P(z)
=
-
\sum_{m\ge 1}
\frac{z^m}{m}
\sum_{k=1}^{N}
a_k^m.
\label{eq:finite_product_logarithmic_expansion}
\end{equation}

In particular,

\begin{equation}
-
\left.
\frac{d}{dz}
\log P(z)
\right|_{z=0}
=
\sum_{k=1}^{N}
a_k.
\label{eq:first_logarithmic_coefficient}
\end{equation}

\end{lemma}

\begin{proof}

Since \(P(z)\) is a finite product,

\begin{equation}
\log P(z)
=
\sum_{k=1}^{N}
\log(1-a_k z).
\label{eq:finite_product_logarithm_sum}
\end{equation}

Using

\begin{equation}
\log(1-w)
=
-
\sum_{m\ge 1}
\frac{w^m}{m}
\label{eq:log_one_minus_w_expansion}
\end{equation}

for \(|w|\) sufficiently small, we obtain

\begin{align}
\log P(z)
&=
-
\sum_{k=1}^{N}
\sum_{m\ge 1}
\frac{(a_k z)^m}{m}
\label{eq:finite_product_expansion_first}
\\
&=
-
\sum_{m\ge 1}
\frac{z^m}{m}
\sum_{k=1}^{N}
a_k^m.
\label{eq:finite_product_expansion_second}
\end{align}

Differentiating at \(z=0\) gives
\eqref{eq:first_logarithmic_coefficient}.

\end{proof}
\subsection[The Phase Product and the Sum SN]{The Phase Product and the Sum \texorpdfstring{\(S_N\)}{SN}}
\label{subsec:phase_product_SN}

We now apply Lemma~\ref{lem:finite_product_logarithmic_coefficients}
to the phases

\begin{equation}
a_k
=
e^{i\gamma_k}.
\label{eq:phase_product_ak}
\end{equation}

\begin{definition}[Phase product]
\label{def:phase_product}

The phase product associated with the first \(N\) ordinates is

\begin{equation}
P_N(z)
=
\prod_{k=1}^{N}
\left(1-z e^{i\gamma_k}\right).
\label{eq:phase_product_definition}
\end{equation}

\end{definition}

The zeros of \(P_N\) are

\begin{equation}
z
=
e^{-i\gamma_k},
\qquad
k=1,\ldots,N.
\label{eq:phase_product_zeros}
\end{equation}

Thus all zeros of \(P_N\) lie on the unit circle.

For each integer \(m\ge 1\), define the \(m\)-th phase sum by

\begin{equation}
S_N^{(m)}
=
\sum_{k=1}^{N}
e^{im\gamma_k}.
\label{eq:higher_phase_sum_definition}
\end{equation}

Then Lemma~\ref{lem:finite_product_logarithmic_coefficients} gives

\begin{equation}
\log P_N(z)
=
-
\sum_{m\ge 1}
\frac{S_N^{(m)}}{m}
z^m.
\label{eq:phase_product_logarithmic_expansion}
\end{equation}

In particular,

\begin{equation}
S_N
=
S_N^{(1)}.
\label{eq:SN_first_phase_mode}
\end{equation}

Equivalently,

\begin{equation}
S_N
=
-
\left.
\frac{d}{dz}
\log P_N(z)
\right|_{z=0}.
\label{eq:SN_logarithmic_derivative}
\end{equation}

Thus \(S_N\) is the first logarithmic coefficient of the phase product
\(P_N\).

\subsection{Polynomial Form and Elementary Symmetric Coefficients}
\label{subsec:elementary_symmetric_coefficients}

It is also useful to consider the monic polynomial

\begin{equation}
Q_N(z)
=
\prod_{k=1}^{N}
\left(z-e^{i\gamma_k}\right).
\label{eq:monic_phase_polynomial}
\end{equation}

Expanding \(Q_N\), one obtains

\begin{equation}
Q_N(z)
=
z^N
-
e_1^{(N)}z^{N-1}
+
e_2^{(N)}z^{N-2}
-
\cdots
+
(-1)^N e_N^{(N)}.
\label{eq:monic_phase_polynomial_expansion}
\end{equation}

Here

\begin{equation}
e_r^{(N)}
=
\sum_{1\le k_1<\cdots<k_r\le N}
e^{i(\gamma_{k_1}+\cdots+\gamma_{k_r})}.
\label{eq:elementary_symmetric_phase_coefficients}
\end{equation}

In particular,

\begin{equation}
e_1^{(N)}
=
\sum_{k=1}^{N}
e^{i\gamma_k}
=
S_N.
\label{eq:first_elementary_coefficient_SN}
\end{equation}

Hence

\begin{equation}
Q_N(z)
=
z^N
-
S_N z^{N-1}
+
e_2^{(N)}z^{N-2}
-
\cdots
+
(-1)^N e_N^{(N)}.
\label{eq:monic_phase_polynomial_SN_coefficient}
\end{equation}

This identifies \(S_N\) as the first elementary symmetric coefficient
of the phases \(e^{i\gamma_k}\).

\subsection{Product Form of the Centroid}
\label{subsec:product_form_centroid}

The centroid was defined by

\begin{equation}
C_N
=
\frac{1}{N}
\sum_{n=1}^{N}
S_n.
\label{eq:CN_definition_product_section}
\end{equation}

For each \(n\le N\), define the partial phase product

\begin{equation}
P_n(z)
=
\prod_{k=1}^{n}
\left(1-z e^{i\gamma_k}\right).
\label{eq:partial_phase_product_definition}
\end{equation}

By \eqref{eq:SN_logarithmic_derivative}, applied with \(N\) replaced by
\(n\), one has

\begin{equation}
S_n
=
-
\left.
\frac{d}{dz}
\log P_n(z)
\right|_{z=0}.
\label{eq:Sn_partial_product_logarithmic_derivative}
\end{equation}

Substitution into \eqref{eq:CN_definition_product_section} gives

\begin{equation}
C_N
=
-
\frac{1}{N}
\sum_{n=1}^{N}
\left.
\frac{d}{dz}
\log P_n(z)
\right|_{z=0}.
\label{eq:CN_sum_logarithmic_derivatives}
\end{equation}

By linearity of differentiation,

\begin{equation}
C_N
=
-
\left.
\frac{d}{dz}
\left(
\frac{1}{N}
\sum_{n=1}^{N}
\log P_n(z)
\right)
\right|_{z=0}.
\label{eq:CN_average_logarithmic_derivative}
\end{equation}
\begin{definition}[Accumulated phase product]
\label{def:accumulated_phase_product}

The accumulated phase product is defined by

\begin{equation}
\mathcal{P}_N(z)
=
\prod_{n=1}^{N}
P_n(z).
\label{eq:accumulated_phase_product_definition}
\end{equation}

\end{definition}

Since

\begin{equation}
\log \mathcal{P}_N(z)
=
\sum_{n=1}^{N}
\log P_n(z),
\label{eq:accumulated_phase_product_logarithm}
\end{equation}

we obtain the exact product identity

\begin{equation}
C_N
=
-
\frac{1}{N}
\left.
\frac{d}{dz}
\log \mathcal{P}_N(z)
\right|_{z=0}.
\label{eq:CN_accumulated_product_logarithmic_derivative}
\end{equation}

Thus \(C_N\) is the normalized first logarithmic coefficient of the
accumulated phase product \(\mathcal{P}_N\).

\subsection{The Triangular Product}
\label{subsec:triangular_phase_product}

Using the definition of \(P_n\),

\begin{equation}
\mathcal{P}_N(z)
=
\prod_{n=1}^{N}
\prod_{k=1}^{n}
\left(1-z e^{i\gamma_k}\right).
\label{eq:accumulated_phase_product_double_product}
\end{equation}

Changing the order of multiplication gives

\begin{equation}
\mathcal{P}_N(z)
=
\prod_{k=1}^{N}
\left(1-z e^{i\gamma_k}\right)^{N-k+1}.
\label{eq:triangular_phase_product}
\end{equation}

Combining
\eqref{eq:CN_accumulated_product_logarithmic_derivative}
and
\eqref{eq:triangular_phase_product},
one obtains

\begin{equation}
C_N
=
-
\frac{1}{N}
\left.
\frac{d}{dz}
\log
\prod_{k=1}^{N}
\left(1-z e^{i\gamma_k}\right)^{N-k+1}
\right|_{z=0}.
\label{eq:CN_triangular_product_logarithmic_derivative}
\end{equation}

Differentiating explicitly gives

\begin{equation}
C_N
=
\frac{1}{N}
\sum_{k=1}^{N}
(N-k+1)e^{i\gamma_k}.
\label{eq:CN_triangular_weighted_sum}
\end{equation}

This is the triangular weighted form of the centroid.

\subsection{Triangular Harmonic Interpretation}
\label{subsec:triangular_harmonic_interpretation}

The usual empirical measure of the phases is

\begin{equation}
\nu_N
=
\frac{1}{N}
\sum_{k=1}^{N}
\delta_{e^{i\gamma_k}}.
\label{eq:uniform_empirical_phase_measure}
\end{equation}

Its first moment is

\begin{equation}
\int_{\mathbb T}
z\,d\nu_N(z)
=
\frac{1}{N}
\sum_{k=1}^{N}
e^{i\gamma_k}
=
\frac{S_N}{N}.
\label{eq:uniform_empirical_phase_first_moment}
\end{equation}

The centroid \(C_N\), by contrast, corresponds to the triangularly
weighted finite measure

\begin{equation}
\mu_N^{\triangle}
=
\frac{1}{N}
\sum_{k=1}^{N}
(N-k+1)
\delta_{e^{i\gamma_k}}.
\label{eq:triangular_phase_measure}
\end{equation}

Its first moment is

\begin{equation}
\int_{\mathbb T}
z\,d\mu_N^{\triangle}(z)
=
\frac{1}{N}
\sum_{k=1}^{N}
(N-k+1)e^{i\gamma_k}
=
C_N.
\label{eq:triangular_phase_measure_first_moment}
\end{equation}

Therefore \(C_N\) is the first harmonic of a triangularly weighted phase
measure, whereas \(S_N/N\) is the first harmonic of the uniformly
weighted empirical phase measure.

\subsection{Local Analytic Interpretation at the Origin}
\label{subsec:local_analytic_interpretation_origin}

From \eqref{eq:phase_product_logarithmic_expansion},

\begin{equation}
\log P_N(z)
=
-
S_N z
+
O(z^2)
\qquad
(z\to 0).
\label{eq:phase_product_local_expansion}
\end{equation}

Thus

\begin{equation}
S_N
=
-
\left.
\frac{d}{dz}
\log P_N(z)
\right|_{z=0}.
\label{eq:SN_local_analytic_interpretation}
\end{equation}

Similarly, \eqref{eq:CN_accumulated_product_logarithmic_derivative}
identifies \(C_N\) as the normalized local logarithmic derivative of
the accumulated phase product \(\mathcal{P}_N\) at the origin.

Consequently, \(S_N\) and \(C_N\) are local logarithmic derivatives at
the origin of two finite products associated with the phases
\(e^{i\gamma_k}\): the ordinary phase product \(P_N\) and the accumulated
phase product \(\mathcal{P}_N\), respectively.

\begin{remark}
\label{rem:product_form_summary}

The identities above are exact finite-\(N\) identities.

They do not depend on the Riemann Hypothesis, on random matrix models,
or on probabilistic assumptions about the zeros.

They show that \(S_N\) is the first logarithmic coefficient of the
ordinary phase product \(P_N\), while \(C_N\) is the normalized first
logarithmic coefficient of the accumulated product
\(\mathcal{P}_N\).

\end{remark}


\section[Supplementary Non-Resonant Triangular Landau--Gonek Structure of SN and CN]{Supplementary Non-Resonant Triangular Landau--Gonek Structure of \texorpdfstring{\(S_N\)}{SN} and \texorpdfstring{\(C_N\)}{CN}}
\label{sec:non_resonant_triangular_lg_structure}

This supplementary section records rigorous structural identities for
the ordinal sum \(S_N\), the height-truncated sums \(S_X(T)\), and the
centroid \(C_N\). The Landau--Gonek phase formulas recalled here are literature-based; the
terminology describing \(C_N\) as a triangular non-resonant transform is
an original structural interpretation of the identities below.

\subsection{Ordinal and height truncations}
\label{subsec:ordinal_and_height_truncations}

The sum \(S_N\) is an ordinal truncation: \(N\) is a positive integer
and counts the first \(N\) ordinates in the ordered list
\eqref{eq:gamma_ordering}. By contrast, height truncations are cut off by
an analytic height \(T\) on the ordinate axis. The counting function
\(N_\zeta(T)\), defined in \eqref{eq:zeta_zero_counting_function},
translates the height \(T\) into the corresponding ordinal number of
zeros. Thus \(N\) is not identified with \(T\).

\begin{definition}
\label{def:continuous_height_phase_sum}

For \(X\in\mathbb{R}_{>0}\) and \(T>0\), define

\begin{equation}
S_X(T)
=
\sum_{0<\gamma\leq T}X^{i\gamma},
\label{eq:continuous_height_phase_sum}
\end{equation}

where zeros are counted with multiplicity and

\begin{equation}
X^{i\gamma}:=\exp(i\gamma\log X)
\label{eq:continuous_power_real_log}
\end{equation}

with the real logarithm of \(X\).

\end{definition}

\begin{theorem}[Height-to-ordinal bridge]
\label{thm:height_to_ordinal_bridge}

For every \(T>0\),

\begin{equation}
S_e(T)
=
\sum_{0<\gamma\leq T}e^{i\gamma}
=
\sum_{k=1}^{N_\zeta(T)}e^{i\gamma_k}
=
S_{N_\zeta(T)}.
\label{eq:height_to_ordinal_bridge}
\end{equation}

\end{theorem}

\begin{proof}
The ordered sequence \eqref{eq:gamma_ordering} lists the positive
ordinates of the nontrivial zeros with multiplicity. By definition,
\(N_\zeta(T)\) is exactly the number of entries in this list whose
ordinate is at most \(T\). Therefore the height-truncated multiset
\(\{\gamma:0<\gamma\leq T\}\), counted with multiplicity, is the same
multiset as \(\{\gamma_k:1\leq k\leq N_\zeta(T)\}\). Summing
\(e^{i\gamma}\) over these two identical multisets gives the first two
equalities in \eqref{eq:height_to_ordinal_bridge}. The last equality is
precisely the ordinal definition \eqref{eq:SN_definition} with
\(N=N_\zeta(T)\).
\end{proof}

\begin{corollary}[Ordinal recovery away from ambiguity]
\label{cor:ordinal_recovery_from_height}

For every height \(T>0\) such that \(N_\zeta(T)=N\),

\begin{equation}
S_N=S_e(T).
\label{eq:ordinal_recovery_from_height}
\end{equation}

\end{corollary}

\begin{proof}
Substitute the hypothesis \(N_\zeta(T)=N\) into
\eqref{eq:height_to_ordinal_bridge}.
\end{proof}

\subsection{Continuous phase sums}
\label{subsec:continuous_phase_sums}

The parameter \(X\) in \eqref{eq:continuous_height_phase_sum} is a
positive real continuous frequency. The definition
\eqref{eq:continuous_power_real_log} avoids any multivalued complex
logarithm.

\begin{lemma}[The continuous frequency \(X=e\)]
\label{lem:e_continuous_frequency_phase}

At the continuous frequency \(X=e\), the height-truncated phase sum is

\begin{equation}
S_e(T)
=
\sum_{0<\gamma\leq T}e^{i\gamma}.
\label{eq:e_continuous_frequency_phase_sum}
\end{equation}

\end{lemma}

\begin{proof}
By the definition of the continuous power with positive real base, for every
\(X\in\mathbb{R}_{>0}\), the identity
\begin{equation}
X^{i\gamma}
=
\exp(i\gamma\log X).
\label{eq:continuous_power_identity_in_e_lemma}
\end{equation}
holds. Taking \(X=e\), and using \(\log e=1\), gives
\begin{equation}
X^{i\gamma}\big|_{X=e}
=
e^{i\gamma}
=
\exp(i\gamma).
\label{eq:e_frequency_power_evaluation}
\end{equation}
Substituting this identity into the definition of \(S_X(T)\) with \(X=e\)
proves the claim.
\end{proof}

\begin{remark}
\label{rem:continuous_not_arithmetic_frequency}

The frequency \(X=e\) is a fixed positive real frequency. In the
arithmetic Landau--Gonek formula, resonances occur at integral prime
powers. With the extended convention

\begin{equation}
\Lambda_{\mathbb R}(X)=0
\end{equation}

away from prime powers, one has

\begin{equation}
\Lambda_{\mathbb R}(e)=0.
\end{equation}

Equivalently, \(e\) is non-resonant because it is not an integral prime
power.
\end{remark}

\subsection{Arithmetic Landau--Gonek resonances}
\label{subsec:arithmetic_landau_gonek_resonances}

Landau--Gonek resonances are arithmetic. In this subsection
\(X\in\mathbb{R}_{>0}\) denotes a continuous frequency, while
\(x\in\mathbb{N}\), \(x>1\), denotes an arithmetic frequency.

\begin{definition}
\label{def:prime_power_resonance_set}

The arithmetic resonance set is

\begin{equation}
\mathcal R
=
\{p^m:\ p\ \text{prime},\ m\in\mathbb N\}.
\label{eq:prime_power_resonance_set}
\end{equation}

\end{definition}

For fixed \(x\in\mathbb N\), \(x>1\), the Landau--Gonek formula under RH
has the phase form

\begin{equation}
\sum_{0<\gamma\leq T}x^{i\gamma}
=
-
\frac{T}{2\pi}
\frac{\Lambda(x)}{\sqrt{x}}
+
O(\log T).
\label{eq:landau_gonek_new_section_phase_form}
\end{equation}

The main term is supported on the prime-power set \(\mathcal R\).

\begin{theorem}[Non-resonance of \(e\)]
\label{thm:e_non_resonant}

The continuous frequency \(X=e\) is not an arithmetic Landau--Gonek
resonance:

\begin{equation}
e\notin\mathcal R.
\label{eq:e_not_in_resonance_set}
\end{equation}

\end{theorem}

\begin{proof}
Every element of \(\mathcal R\) is a positive integer, so
\(\mathcal R\subset\mathbb N\). The number \(e\) is irrational, hence
\(e\notin\mathbb N\). Therefore \(e\notin\mathcal R\).
\end{proof}

\begin{remark}
\label{rem:no_mangoldt_at_e}

In the extended real-frequency notation,
\(\Lambda_{\mathbb R}(X)=0\) away from integral prime powers. Thus
\(\Lambda_{\mathbb R}(e)=0\), equivalently \(X=e\) is non-resonant
because it is not an integral prime power.
\end{remark}

\subsection{The centroid as a triangular non-resonant transform}
\label{subsec:centroid_triangular_non_resonant_transform}

The centroid identity is purely ordinal and contains no height variable.
It follows from the definitions in Section~\ref{sec:SN_definition}.

\begin{theorem}[Triangular centroid identity]
\label{thm:triangular_centroid_identity}

For every positive integer \(N\),

\begin{equation}
C_N
=
\frac{1}{N}\sum_{n=1}^{N}S_n
=
\frac{1}{N}\sum_{n=1}^{N}\sum_{k=1}^{n}e^{i\gamma_k}
=
\frac{1}{N}\sum_{k=1}^{N}(N-k+1)e^{i\gamma_k}.
\label{eq:triangular_centroid_identity}
\end{equation}

\end{theorem}

\begin{proof}
The first equality is the definition \eqref{eq:CN_definition}. Replacing
\(S_n\) by its defining sum gives the second equality, which is
\eqref{eq:CN_double_sum}. In the double sum, the term
\(e^{i\gamma_k}\) appears exactly for those indices \(n\) with
\(k\leq n\leq N\). There are \(N-k+1\) such indices. Collecting equal
terms gives the final expression, already recorded in
\eqref{eq:CN_weighted}.
\end{proof}

\begin{corollary}[Triangular non-resonant centroid transform]

The centroid \(C_N\) is a triangularly weighted non-resonant phase sum associated with the continuous frequency \(X=e\).

\end{corollary}

\begin{proof}

By the triangular centroid identity,
\begin{equation}
C_N
=
\frac{1}{N}
\sum_{k=1}^{N}
(N-k+1)e^{i\gamma_k}.
\label{eq:triangular_centroid_identity_recalled}
\end{equation}
Thus \(C_N\) is a triangularly weighted sum of the phases \(e^{i\gamma_k}\).

By the continuous-frequency identity at \(X=e\), these phases coincide with the continuous-frequency phases \(X^{i\gamma_k}\) evaluated at \(X=e\). Since \(e\) is not contained in the arithmetic resonance set
\begin{equation}
\mathcal R
=
\{p^m:\ p\ \text{prime},\ m\in\mathbb N\},
\label{eq:prime_power_resonance_set_recalled}
\end{equation}
the triangularly weighted ordinal sum defining \(C_N\) is non-resonant in the Landau--Gonek prime-power sense.

\end{proof}

\subsection{Equidistribution consequences}
\label{subsec:equidistribution_consequences_non_resonant}

Hlawka's theorem, applied with \(\alpha=1/(2\pi)\), implies the
normalized cancellation already stated in \eqref{eq:SN_normalized_cancellation}.
The next consequence records the same fact in the notation of this
section.

\begin{corollary}[Normalized cancellation at \(X=e\)]
\label{cor:non_resonant_normalized_cancellation}

The ordinal phase sums at the continuous frequency \(X=e\) satisfy

\begin{equation}
\frac{S_N}{N}\longrightarrow 0
\qquad
(N\to\infty).
\label{eq:non_resonant_normalized_cancellation}
\end{equation}

\end{corollary}

\begin{proof}
Equation \eqref{eq:SN_normalized_cancellation} is obtained from Hlawka's
equidistribution theorem and Weyl's criterion with
\(\alpha=1/(2\pi)\). Since \(e^{i\gamma_k}\) is the phase associated with
\(X=e\) by Lemma~\ref{lem:e_continuous_frequency_phase}, the result is
exactly \eqref{eq:non_resonant_normalized_cancellation}.
\end{proof}



\begin{theorem}[Sublinear centroid growth under normalized cancellation]
\label{thm:sublinear_centroid_growth}

Assume that

\begin{equation}
\frac{S_N}{N}\longrightarrow 0
\qquad
(N\to\infty).
\label{eq:sublinear_centroid_assumption}
\end{equation}

Then

\begin{equation}
\frac{C_N}{N}\longrightarrow 0
\qquad
(N\to\infty).
\label{eq:sublinear_centroid_conclusion_normalized}
\end{equation}

Equivalently,

\begin{equation}
C_N=o(N).
\label{eq:sublinear_centroid_conclusion_little_o}
\end{equation}

\end{theorem}

\begin{proof}
By definition,

\begin{equation}
C_N
=
\frac{1}{N}\sum_{n=1}^{N}S_n.
\label{eq:sublinear_centroid_definition_used}
\end{equation}

Therefore

\begin{equation}
\frac{C_N}{N}
=
\frac{1}{N^2}\sum_{n=1}^{N}S_n.
\label{eq:centroid_over_N_average_partial_sums}
\end{equation}

Let \(\varepsilon>0\). Since

\begin{equation}
\frac{S_n}{n}\longrightarrow 0,
\label{eq:partial_sum_normalized_n_limit}
\end{equation}

there exists \(N_0\in\mathbb{N}\) such that

\begin{equation}
|S_n|\leq \varepsilon n
\label{eq:partial_sum_epsilon_bound}
\end{equation}

for every \(n\geq N_0\). For \(N\geq N_0\), split the sum as

\begin{equation}
\frac{C_N}{N}
=
\frac{1}{N^2}\sum_{n=1}^{N_0-1}S_n
+
\frac{1}{N^2}\sum_{n=N_0}^{N}S_n.
\label{eq:centroid_sum_split}
\end{equation}

Taking absolute values gives

\begin{equation}
\left|\frac{C_N}{N}\right|
\leq
\frac{1}{N^2}\sum_{n=1}^{N_0-1}|S_n|
+
\frac{1}{N^2}\sum_{n=N_0}^{N}|S_n|.
\label{eq:centroid_sum_split_abs}
\end{equation}

The first numerator is fixed, hence

\begin{equation}
\frac{1}{N^2}\sum_{n=1}^{N_0-1}|S_n|
\longrightarrow 0.
\label{eq:centroid_initial_block_vanishes}
\end{equation}

For the second term, using \eqref{eq:partial_sum_epsilon_bound}, we
obtain

\begin{equation}
\frac{1}{N^2}\sum_{n=N_0}^{N}|S_n|
\leq
\frac{\varepsilon}{N^2}\sum_{n=N_0}^{N}n
\leq
\frac{\varepsilon}{N^2}\sum_{n=1}^{N}n.
\label{eq:centroid_tail_bound}
\end{equation}

Since

\begin{equation}
\sum_{n=1}^{N}n
=
\frac{N(N+1)}{2},
\label{eq:sum_first_N_integers_centroid}
\end{equation}

it follows that

\begin{equation}
\frac{\varepsilon}{N^2}\sum_{n=1}^{N}n
=
\frac{\varepsilon}{2}\left(1+\frac{1}{N}\right).
\label{eq:centroid_tail_bound_evaluated}
\end{equation}

Therefore

\begin{equation}
\limsup_{N\to\infty}
\left|\frac{C_N}{N}\right|
\leq
\frac{\varepsilon}{2}.
\label{eq:centroid_limsup_epsilon}
\end{equation}

Since \(\varepsilon>0\) is arbitrary, this proves
\eqref{eq:sublinear_centroid_conclusion_normalized}, equivalently
\eqref{eq:sublinear_centroid_conclusion_little_o}.

\end{proof}

\begin{corollary}[Equidistribution does not force centroid convergence]
\label{cor:equidistribution_does_not_force_CN_zero}

The normalized cancellation in
Corollary~\ref{cor:non_resonant_normalized_cancellation} implies only

\begin{equation}
C_N=o(N),
\label{eq:centroid_little_o_only_statement}
\end{equation}

and does not imply

\begin{equation}
C_N\longrightarrow 0.
\label{eq:centroid_zero_not_implied_statement}
\end{equation}

\end{corollary}

\begin{proof}
The sublinear estimate follows from
Theorem~\ref{thm:sublinear_centroid_growth}. However, sublinear growth
alone does not imply convergence to zero. For example, the scalar sequence

\begin{equation}
a_N=\sqrt N
\label{eq:scalar_sublinear_example}
\end{equation}

satisfies

\begin{equation}
\frac{a_N}{N}\longrightarrow 0,
\label{eq:scalar_sublinear_example_normalized}
\end{equation}

but does not converge to zero. Therefore normalized cancellation of the
endpoint does not, by itself, force centroid convergence to zero.
\end{proof}

\subsection{Structural conclusion}
\label{subsec:structural_conclusion_non_resonant}

The preceding results separate the ordinal, height-truncated, and
arithmetic layers of the construction. The sum \(S_N\) is an ordinal
phase walk, while \(S_X(T)\) is a height-truncated phase sum; the counting
function \(N_\zeta(T)\) provides the bridge between the two. The phase
\(e^{i\gamma}\) corresponds to the continuous frequency \(X=e\). In the
arithmetic Landau--Gonek formula the resonant frequencies are the
integral prime powers. The present use of \(X=e\) is a fixed positive
real frequency.

The Chebyshev and logarithmic-derivative viewpoint explains the
arithmetic origin of the resonant frequencies: the weights \(\Lambda(x)\)
arise from

\begin{equation}
-\frac{\zeta'(s)}{\zeta(s)}
=
\sum_{\ell=1}^{\infty}\frac{\Lambda(\ell)}{\ell^s},
\label{eq:structural_conclusion_log_derivative_bridge}
\end{equation}

and the Landau--Gonek main terms are supported on prime powers. With the
extended real convention \(\Lambda_{\mathbb R}(X)=0\) away from prime
powers, one has \(\Lambda_{\mathbb R}(e)=0\). Equivalently, \(e\) is
non-resonant because it is not an integral prime power. Thus neither the
phase walk \(S_N\) nor the triangular centroid \(C_N\) receives a direct
Chebyshev--Landau main term. The centroid is a triangular non-resonant
phase transform at the continuous frequency \(X=e\).

This interpretation is compatible with Weyl and Hlawka
equidistribution and with the Landau--Gonek prime-power resonance
formula. Equidistribution gives normalized cancellation of the
increments, and Theorem~\ref{thm:sublinear_centroid_growth} gives the
sublinear bound \(C_N=o(N)\). Landau--Gonek identifies arithmetic
resonances at prime powers, whereas the present construction uses the
non-arithmetic frequency \(X=e\). The conclusion is structural: it
separates the ordinal and arithmetic layers and records the
non-resonant role of the triangular centroid.

\section{Conclusions}
\label{sec:conclusions}

The paper isolates the following outputs.

\subsection{Complete-zero logarithmic confinement}
\label{subsec:conclusion_complete_zero_confinement}
\leavevmode\par\nobreak

Landau's fixed-real formula at \(X=e\) gives a logarithmic confinement
baseline for the complete-zero walk:

\begin{equation}
Z_N=O(\log N),
\qquad
\mathcal C_N=O(\log N),
\qquad
\mathcal M_0(N)=O\!\left((\log N)^2\right).
\label{eq:conclusion_complete_zero_bounds}
\end{equation}

Under RH this specializes to the original unit phase walk through

\begin{equation}
Z_N=S_N,
\qquad
\mathcal C_N=C_N,
\qquad
\mathcal M_0(N)=M_0(N).
\label{eq:conclusion_RH_specialized_objects}
\end{equation}

\subsection{Exact centered pair statistic}
\label{subsec:conclusion_exact_pair_statistic}
\leavevmode\par\nobreak

In the complete-zero setting,

\begin{equation}
\mathcal M_0(N)
=
\sum_{j,k=1}^{N}
W_N(j,k)
e^{\rho_j-\frac12}
e^{\overline{\rho_k}-\frac12}.
\label{eq:conclusion_complete_zero_pair_statistic}
\end{equation}

Under RH, this becomes the exact centered pair statistic

\begin{equation}
M_0(N)
=
\sum_{j,k=1}^{N}
W_N(j,k)e^{i(\gamma_j-\gamma_k)}
\label{eq:conclusion_RH_exact_pair_statistic}
\end{equation}

Here \(W_N\) is the exact centering kernel

\begin{equation}
W_N(j,k)
=
\frac{N-\max(j,k)+1}{N}
-
\frac{(N-j+1)(N-k+1)}{N^2}.
\label{eq:conclusion_WN_definition}
\end{equation}

It is the weight forced by subtracting the triangular centroid before
forming the quadratic moment.

\subsection{Brownian-bridge/Green kernel}
\label{subsec:conclusion_brownian_green_kernel}
\leavevmode\par\nobreak

After reversing indices, the kernel becomes

\begin{equation}
K_N(r,s)
=
\frac{\min(r,s)}{N}
-
\frac{rs}{N^2},
\label{eq:conclusion_brownian_green_kernel}
\end{equation}

the discrete Brownian-bridge covariance kernel, equivalently the Green
kernel of the discrete Dirichlet Laplacian. This is the central
structural point of the paper.

\subsection{Lifted phases and Green-modal coordinates}
\label{subsec:conclusion_lifted_phases_green_modal_coordinates}
\leavevmode\par\nobreak

There is no modal shortcut hidden in the notation. The lifted phase
representation in Lemma~\ref{lem:lifted_phase_representation} separates
the angular representative \(\theta_k\) from the winding number \(B_k\).
For \(t=1\), the winding is invisible in \(e^{i\gamma_k}\); for real
frequency deformations \(e^{it\gamma_k}\), and in particular in the
Montgomery scaling, the winding data are part of the phase.

The modal basis is the normalized Dirichlet sine basis
\eqref{eq:dirichlet_sine_modes}. Its orthonormality is proved in
Lemma~\ref{lem:dirichlet_sine_orthonormality}, and exact modal inversion
is Lemma~\ref{lem:exact_modal_inversion}. Thus the coefficient vector
\(A(N;t)\) is genuine coordinate data, not a heuristic compression.

With these coordinates, the Green kernel diagonalizes and yields

\begin{equation}
\mathcal E_N(t)
=
A(N;t)^*\Lambda_NA(N;t)
=
\sum_{m=1}^{N-1}
\lambda_m|A_m(N;t)|^2.
\label{eq:conclusion_green_energy_modal_decomposition}
\end{equation}

Under RH at \(t=1\), this is the centered phase-walk moment \(M_0(N)\).

\subsection{Exact recovery versus scalar energy bounds}
\label{subsec:conclusion_exact_recovery_scalar_bounds}
\leavevmode\par\nobreak

Put

\begin{equation}
S_N(t)
=
\sum_{k=1}^{N}e^{it\gamma_k}.
\label{eq:conclusion_t_deformed_endpoint}
\end{equation}

The endpoint is recovered exactly from the Green-modal coefficient data:

\begin{equation}
S_N(t)
=
e^{it\gamma_1}
+
\sum_{\substack{1\le m\le N-1\\ m\ \mathrm{odd}}}
u_mA_m(N;t).
\label{eq:conclusion_endpoint_projection}
\end{equation}

Even modes vanish from the endpoint sum only because their projection
onto the constant vector is zero; they still contribute to the centered
Green energy.

The scalar Green energy alone does not recover the endpoint or the
centroid: it discards modal direction and phase information. Exact
recovery uses the coefficient vector \(A(N;t)\), equivalently the
weighted modal data \(\Lambda_NA(N;t)\), as in
\eqref{eq:green_metric_endpoint_recovery_identity}. The scalar energy
gives only one-sided deweighting estimates. From
\eqref{eq:finite_green_deweighting} and
\eqref{eq:centroid_green_deweighting_bound},

\begin{equation}
\left|S_N(t)-e^{it\gamma_1}\right|^2
\le
2N\,\mathcal E_N(t),
\label{eq:conclusion_endpoint_weak_green_bound}
\end{equation}

and

\begin{equation}
\left|C_N(t)-e^{it\gamma_1}\right|^2
\le
(N-1)\mathcal E_N(t).
\label{eq:conclusion_centroid_weak_green_bound}
\end{equation}

Consequently, if \(\mathcal E_N(1)=O((\log N)^2)\), the Green bounds give
only the weaker estimates

\begin{equation}
S_N=O(\sqrt N\log N),
\qquad
C_N=O(\sqrt N\log N).
\label{eq:conclusion_weak_green_bounds_at_t_one}
\end{equation}

The logarithmic bounds in \eqref{eq:conclusion_complete_zero_bounds} are
stronger and come from Landau's fixed-real formula at \(X=e\), not from
the scalar Green-energy inequality.

\subsection{Green-modal Montgomery formulation}
\label{subsec:conclusion_green_modal_montgomery}
\leavevmode\par\nobreak

The smoothed Montgomery form factor can be written as

\begin{equation}
F_N(\tau)
=
\frac1N
\int_{-\infty}^{\infty}
e^{-2|y|}
\left|
e^{i(\tau\log(T_N/2\pi)+y)\gamma_1}
+
\sum_{\substack{1\le m\le N-1\\ m\ \mathrm{odd}}}
u_mA_m(N;\tau\log(T_N/2\pi)+y)
\right|^2dy.
\label{eq:conclusion_green_modal_montgomery}
\end{equation}

Thus the Montgomery--GUE form-factor target can be expressed as a
statement about the Green-modal endpoint projection:

\begin{equation}
F_N(\tau)\to\min(|\tau|,1).
\label{eq:conclusion_montgomery_target}
\end{equation}

The finite modal bounds on individual \(A_m(N;t)\) are unconditional, but
they do not determine this limit. The conditional modal covariance
criterion is stated in
Corollary~\ref{cor:modal_covariance_montgomery_criterion}: the limiting
form factor is governed by the rank-one contraction
\(N^{-1}u^*\Sigma_N(\tau)u\), not by pointwise modal bounds alone.

The final logical division is therefore sharp: Landau gives the strong
logarithmic confinement estimates; the Green kernel gives the exact
structural identity, the modal decomposition, and the exact deweighting
of linear functionals. The scalar Green energy supplies canonical upper
bounds, not exact recovery of endpoint-type quantities.

\subsection{RVM residual decomposition}
\label{subsec:conclusion_rvm_residual_decomposition}
\leavevmode\par\nobreak

Here RVM means Riemann--von Mangoldt. Let
\(\widetilde\gamma_k=N_0^{-1}(k)\) denote the synthetic ordinates
obtained by inverting the smooth Riemann--von Mangoldt counting main
term. These ordinates form the rigid RVM background against which the
true zeta ordinates \(\gamma_k\) are compared, and write
\(d_k=\gamma_k-\widetilde\gamma_k\) for the residual displacement.

The modes split as

\begin{equation}
A_m(N;t)=A_m^{\mathrm{RVM}}(N;t)+R_m(N;t),
\label{eq:conclusion_rvm_modal_split}
\end{equation}

with

\begin{equation}
\sum_{m=1}^{N-1}|R_m(N;t)|^2
=
\sum_{k=2}^{N}|e^{itd_k}-1|^2
\le
4(N-1).
\label{eq:conclusion_rvm_residual_bound}
\end{equation}

Equivalently,

\begin{equation}
\frac1N
\sum_{m=1}^{N-1}|R_m(N;t)|^2
\le
4.
\label{eq:conclusion_rvm_residual_normalized_bound}
\end{equation}

Thus \(A_m^{\mathrm{RVM}}(N;t)\) is the modal coefficient computed from
the synthetic RVM ordinates \(\widetilde\gamma_k\), while \(R_m(N;t)\) is
the residual modal difference contributed by replacing those synthetic
ordinates with the true zero ordinates. In other words, \(R_m\) measures
the mode-by-mode remainder between the synthetic RVM grid and the actual
zeros.

The main contribution of the paper is the identification of the
Brownian-bridge/Green kernel hidden in the centered phase-walk geometry.
This kernel converts the centroid problem into an exact weighted
pair-correlation statistic and provides a natural Green-modal language
in which the Montgomery--GUE form factor can be expressed.
The centered moment gives an exact Green-kernel weighted pair statistic.
The Green-modal coordinates give an exact endpoint representation of the
Montgomery form factor, while the scalar Green energy gives the
corresponding one-sided deweighting estimate.

\subsection{Structural separation}
\label{subsec:conclusion_structural_separation}
\leavevmode\par\nobreak

The fixed-frequency objects \(S_N(1)\), \(C_N\), and \(M_0(N)\)
describe non-resonant Landau cancellation and Green-weighted centered
geometry. The Montgomery--GUE statistics are seen through the deformed
endpoint family \(S_N(t)\). In Fourier variables this gives the form
factor, while in the unfolded pair variable

\begin{equation}
\xi_{jk}
=
\frac{\log(T_0/2\pi)}{2\pi}
(\gamma_j-\gamma_k)
\label{eq:conclusion_unfolded_pair_variable_xi}
\end{equation}

it gives the sine-kernel pair correlation. The Green-modal vector
\(A(N;t)\) is the common coordinate system linking these two regimes.

\begin{thebibliography}{99}
\sloppy
\hbadness=10000

\bibitem{Riemann1859}
B. Riemann,
\emph{\"Uber die Anzahl der Primzahlen unter einer gegebenen Gr\"osse},
Monatsberichte der Berliner Akademie,
1859,
671--680.

Reprinted in:

B. Riemann,
\emph{Gesammelte Mathematische Werke},
Dover Publications,
1953.

\bibitem{Titchmarsh1986}
E. C. Titchmarsh,
\emph{The Theory of the Riemann Zeta-Function},
Second Edition,
revised by D. R. Heath-Brown,
Oxford University Press,
1986.

\bibitem{Edwards2001}
H. M. Edwards,
\emph{Riemann's Zeta Function},
Dover Publications,
2001.

\bibitem{IwaniecKowalski2004}
H. Iwaniec and E. Kowalski,
\emph{Analytic Number Theory},
American Mathematical Society,
2004.

\bibitem{Weyl1916}
H. Weyl,
\emph{\"Uber die Gleichverteilung von Zahlen mod. Eins},
Mathematische Annalen
\textbf{77} (1916),
313--352.
doi:10.1007/BF01475864.

\bibitem{Rademacher1956}
H. Rademacher,
\emph{On the distribution of the values of the zeta function},
Colloquium Mathematicum
\textbf{4} (1956),
79--81.

\bibitem{Hlawka1975}
E. Hlawka,
\emph{\"Uber die Gleichverteilung gewisser Folgen,
die mit den Nullstellen der Zetafunktion zusammenh\"angen},
Sitzungsberichte der
\"Osterreichischen Akademie der Wissenschaften,
Mathematisch-Naturwissenschaftliche Klasse,
Abteilung II
\textbf{184} (1975),
459--471.

\bibitem{Landau1911}
E. Landau,
\emph{\"Uber die Nullstellen der Zetafunktion},
Mathematische Annalen
\textbf{71} (1911),
548--564.
doi:10.1007/BF01456881.

\bibitem{Gonek1985}
S. M. Gonek,
\emph{A formula of Landau and mean values of \(\zeta(s)\)},
in:
\emph{Topics in Analytic Number Theory},
edited by S. W. Graham and J. D. Vaaler,
University of Texas Press,
Austin,
1985,
92--97.

\bibitem{Gonek1993}
S. M. Gonek,
\emph{An explicit formula of Landau and its applications to the theory of the zeta-function},
in:
\emph{A Tribute to Emil Grosswald:
Number Theory and Related Analysis},
Contemporary Mathematics
\textbf{143},
American Mathematical Society,
1993,
395--413.

\bibitem{Odlyzko1987}
A. M. Odlyzko,
\emph{On the distribution of spacings between zeros of the zeta function},
Mathematics of Computation
\textbf{48} (1987),
273--308.
doi:10.1090/S0025-5718-1987-0866115-0.

\bibitem{KaczorowskiLanguascoPerelli2000}
J. Kaczorowski,
A. Languasco,
A. Perelli,
\emph{A note on Landau's formula},
Functiones et Approximatio Commentarii Mathematici
\textbf{28} (2000),
173--186.
doi:10.7169/facm/1538186693.

\bibitem{Montgomery1973}
H. L. Montgomery,
\emph{The pair correlation of zeros of the zeta function},
in:
\emph{Analytic Number Theory},
Proceedings of Symposia in Pure Mathematics
\textbf{24},
American Mathematical Society,
1973,
181--193.

\bibitem{Borwein2008}
P. Borwein,
S. Choi,
B. Rooney,
A. Weirathmueller,
\emph{The Riemann Hypothesis:
A Resource for the Afficionado and Virtuoso Alike},
Springer,
2008.

\bibitem{FordZaharescu2005}
K. Ford and A. Zaharescu,
\emph{Unnormalized Differences Between Zeros of the Riemann Zeta Function},
Acta Arithmetica
\textbf{117} (2005),
163--176.
doi:10.4064/aa117-2-5.

\bibitem{FordSoundararajanZaharescu2024}
K. Ford,
K. Soundararajan,
A. Zaharescu,
\emph{The Distribution of the Ordinates of the Zeros of the Riemann Zeta Function},
arXiv:2407.14173,
2024.

\bibitem{Fujii1978}
A. Fujii,
\emph{On the uniformity of the distribution of the zeros of the
Riemann zeta function},
Journal f\"ur die reine und angewandte Mathematik
\textbf{302} (1978),
167--205.

\end{thebibliography}

\end{document}
